\documentclass[11pt,letterpaper]{article}

% --- Idioma, tipografia y pagina -------------------------------------------
\usepackage[T1]{fontenc}
\usepackage[utf8]{inputenc}
\usepackage[spanish,es-nodecimaldot]{babel}
\usepackage{lmodern}
\usepackage{microtype}
\usepackage[margin=2.4cm]{geometry}
\usepackage{setspace}
\setstretch{1.08}

% --- Color, cajas, tablas y listas -----------------------------------------
\usepackage{xcolor}
\definecolor{AzulParcial}{HTML}{24527A}
\definecolor{AzulClaro}{HTML}{EAF2F8}
\definecolor{GrisTexto}{HTML}{374151}
\definecolor{MoradoCaja}{HTML}{7C3AED}
\definecolor{MoradoClaro}{HTML}{F3E8FF}
\definecolor{NaranjaCaja}{HTML}{EA580C}
\definecolor{NaranjaClaro}{HTML}{FFF7ED}
\definecolor{VerdeCaja}{HTML}{15803D}
\definecolor{VerdeClaro}{HTML}{F0FDF4}
\definecolor{AzulLinea}{HTML}{2563EB}
\definecolor{VerdeLinea}{HTML}{16A34A}
\definecolor{NaranjaLinea}{HTML}{EA580C}
\definecolor{MoradoLinea}{HTML}{7C3AED}
\definecolor{RosaLinea}{HTML}{DB2777}
\definecolor{TurquesaLinea}{HTML}{0F766E}

% --- Colores semánticos para tipos de bloque en diagramas de flujo ---------
\definecolor{BordeInicio}{HTML}{2563EB}
\definecolor{FondoInicio}{HTML}{EFF6FF}
\definecolor{BordeProceso}{HTML}{475569}
\definecolor{FondoProceso}{HTML}{F8FAFC}
\definecolor{BordeEntrada}{HTML}{0D9488}
\definecolor{FondoEntrada}{HTML}{F0FDFA}
\definecolor{BordeDecision}{HTML}{D97706}
\definecolor{FondoDecision}{HTML}{FFFBEB}
\definecolor{BordeSalida}{HTML}{7C3AED}
\definecolor{FondoSalida}{HTML}{FAF5FF}
\definecolor{BordeManual}{HTML}{E11D48}
\definecolor{FondoManual}{HTML}{FFF1F2}
\definecolor{BordeConector}{HTML}{4B5563}
\definecolor{FondoConector}{HTML}{F3F4F6}
\usepackage[most]{tcolorbox}
\usepackage{booktabs}
\usepackage{longtable}
\usepackage{array}
\usepackage{calc}
\newcounter{none}
\setlength{\tabcolsep}{3pt}
\hfuzz=0.5pt
\AtBeginEnvironment{longtable}{\footnotesize}
\usepackage{enumitem}
\usepackage{graphicx}
\usepackage{float}
\usepackage{adjustbox}

% --- Matematicas y codigo ---------------------------------------------------
\usepackage{amsmath,amssymb}
\usepackage{listings}
\lstdefinestyle{apuntes}{
  basicstyle=\ttfamily\small,
  backgroundcolor=\color{black!3},
  frame=single,
  rulecolor=\color{black!15},
  breaklines=true,
  showstringspaces=false,
  columns=fullflexible
}
\lstset{style=apuntes}
\providecommand{\tightlist}{%
  \setlength{\itemsep}{0pt}\setlength{\parskip}{0pt}}
\providecommand{\passthrough}[1]{#1}

% --- Diagramas de flujo ----------------------------------------------------
\usepackage{tikz}
\usetikzlibrary{arrows.meta,calc,positioning,shapes.geometric,shapes.symbols}

\tikzset{
  flujo/.style={
    >=Latex,
    node distance=9mm and 16mm,
    every path/.style={draw=GrisTexto,thick,-{Latex[length=2.4mm]}},
    every node/.style={font=\small,align=center,text=GrisTexto}
  },
  etiqueta/.style={
    draw=none,fill=white,inner sep=1.5pt
  },
  iniciofin/.style={
    ellipse,draw=BordeInicio,very thick,fill=FondoInicio,
    minimum width=31mm,minimum height=10mm
  },
  proceso/.style={
    rectangle,draw=BordeProceso,very thick,fill=FondoProceso,
    rounded corners=1pt,minimum width=30mm,minimum height=10mm,
    inner xsep=4mm
  },
  decision/.style={
    diamond,draw=BordeDecision,very thick,fill=FondoDecision,
    aspect=2,inner sep=1.5pt,minimum width=30mm,minimum height=12mm
  },
  entradaSalida/.style={
    trapezium,trapezium left angle=70,trapezium right angle=110,
    draw=BordeEntrada,very thick,fill=FondoEntrada,
    minimum width=32mm,minimum height=10mm,inner xsep=4mm
  },
  documento/.style={
    tape,tape bend top=none,tape bend bottom=in and out,
    tape bend height=5pt,
    draw=BordeSalida,very thick,fill=FondoSalida,
    minimum width=35mm,minimum height=12mm,
    inner xsep=4mm,inner ysep=2.5mm
  },
  operacionManual/.style={
    trapezium,trapezium angle=110,
    draw=BordeManual,very thick,fill=FondoManual,
    minimum width=42mm,minimum height=12mm
  },
  conector/.style={
    circle,draw=BordeConector,very thick,fill=FondoConector,
    minimum size=7mm,inner sep=0pt
  }
}

% --- Encabezados y enlaces -------------------------------------------------
\usepackage{fancyhdr}
\usepackage[hidelinks]{hyperref}
\pagestyle{fancy}
\fancyhf{}
\lhead{Cuadernillo de actividades}
\rhead{Parcial \NumeroParcial}
\cfoot{\thepage}
\setlength{\headheight}{14pt}

% --- Datos que se modifican una sola vez -----------------------------------
\newcommand{\Asignatura}{Razonamiento computacional y diseño de algoritmos}
\newcommand{\NumeroParcial}{1}
\newcommand{\Docente}{Walter Alexander Mata López}
\newcommand{\Estudiante}{}
\newcommand{\Grupo}{}

% --- Entornos reutilizables ------------------------------------------------
\newcounter{actividad}
\newenvironment{actividad}[1]{%
  \refstepcounter{actividad}
  \begin{tcolorbox}[
    enhanced,breakable,
    colback=VerdeClaro,colframe=VerdeCaja,
    colbacktitle=VerdeCaja,coltitle=white,
    fonttitle=\bfseries,
    title={Actividad \theactividad: #1},
    arc=3mm,boxrule=0.8pt,
    attach boxed title to top left={xshift=5mm,yshift=-2.5mm},
    boxed title style={arc=2mm,boxrule=0pt,left=3mm,right=3mm},
    top=7mm,left=4mm,right=4mm,bottom=4mm
  ]
}{\end{tcolorbox}}

\newtcolorbox{objetivo}{
  enhanced,breakable,
  colback=MoradoClaro,colframe=MoradoCaja,
  boxrule=0pt,leftrule=4pt,arc=2.5mm,
  title=Objetivo,fonttitle=\bfseries,
  colbacktitle=MoradoClaro,coltitle=MoradoCaja,
  titlerule=0pt,left=4mm,right=4mm,bottom=3mm
}

\newcommand{\espaciorespuesta}[1][3cm]{%
  \par\vspace{2mm}
  \begin{tcolorbox}[
    enhanced,colback=NaranjaClaro,colframe=NaranjaCaja,
    colbacktitle=NaranjaCaja,coltitle=white,
    height=#1,title=Respuesta,fonttitle=\bfseries,
    arc=2.5mm,boxrule=0.7pt,
    left=3mm,right=3mm,bottom=3mm
  ]
  \end{tcolorbox}
}

\begin{document}
\shorthandoff{"}
\setcounter{secnumdepth}{0}

\hypersetup{pageanchor=false}
\begin{titlepage}
  \thispagestyle{empty}
  \centering
  \vspace*{2.2cm}
  {\color{AzulParcial}\rule{\textwidth}{1.5pt}}\\[1.1cm]
  {\Huge\bfseries Cuadernillo de actividades\par}
  \vspace{3mm}
  {\LARGE Parcial \NumeroParcial\par}
  \vspace{8mm}
  {\Large \Asignatura\par}
  \vspace{1.1cm}
  {\color{AzulParcial}\rule{\textwidth}{1.5pt}}\\[1.4cm]

  \begin{tabular}{@{}rl@{}}
    \textbf{Estudiante:} & \Estudiante \\[3mm]
    \textbf{Grupo:}      & \Grupo \\[3mm]
    \textbf{Docente:}    & \Docente
  \end{tabular}

  \vfill
  {\large Material de trabajo y evidencias de aprendizaje\par}
\end{titlepage}

\hypersetup{pageanchor=true}
\pagenumbering{Roman}
\tableofcontents
\newpage
\pagenumbering{arabic}

\begin{center}\rule{0.5\linewidth}{0.5pt}\end{center}

\section{1. Introducción}\label{introducciuxf3n}

La inteligencia artificial y la ciencia de datos requieren resolver
problemas complejos, desde procesar grandes volúmenes de información
hasta entrenar modelos de aprendizaje automático. No obstante, el
\textbf{razonamiento computacional} es una disciplina general que
trasciende estas áreas, pues permite descomponer cualquier problema y
formular soluciones estructuradas para que una computadora pueda
ejecutarlas con precisión.

Este proceso se apoya principalmente en dos estrategias:

\begin{enumerate}
\def\labelenumi{\arabic{enumi}.}
\tightlist
\item
  \textbf{Abstracción:} Identificar los aspectos esenciales de una
  situación e ignorar los detalles secundarios para concentrarse en la
  información crítica.
\item
  \textbf{Descomposición:} Dividir un problema complejo en partes más
  pequeñas y manejables que puedan resolverse paso a paso.
\end{enumerate}

\begin{center}\rule{0.5\linewidth}{0.5pt}\end{center}

\section{2. Concepto de algoritmo y
granularidad}\label{concepto-de-algoritmo-y-granularidad}

Un \textbf{algoritmo} es un conjunto ordenado, finito y no ambiguo de
instrucciones lógicas que conducen a la solución de un problema
específico. Para ser formalmente correcto, todo algoritmo debe cumplir
tres propiedades fundamentales:

\begin{itemize}
\tightlist
\item
  \textbf{Preciso:} Cada paso debe estar definido con exactitud, sin
  margen de interpretación.
\item
  \textbf{Determinista:} Al procesar las mismas entradas, debe generar
  invariablemente el mismo resultado.
\item
  \textbf{Finito:} Debe concluir tras un número determinado de
  operaciones.
\end{itemize}

\subsection{2.1 El nivel de
granularidad}\label{el-nivel-de-granularidad}

Al diseñar un algoritmo, el \textbf{nivel de granularidad} representa el
grado de detalle con el que se especifican las instrucciones:

\begin{itemize}
\tightlist
\item
  \textbf{Granularidad baja (macro-pasos):} Describe operaciones
  generales (por ejemplo, \emph{``limpiar la base de datos''}). Es útil
  para conceptualizar la arquitectura global de una solución, aunque una
  computadora no puede ejecutar directamente instrucciones tan amplias.
\item
  \textbf{Granularidad alta (micro-pasos):} Expresa operaciones
  elementales y atómicas, al nivel de instrucciones aritméticas y
  lógicas que el procesador interpreta directamente.
\end{itemize}

\subsection{2.2 Algoritmos en la vida
cotidiana}\label{algoritmos-en-la-vida-cotidiana}

Antes de codificar en un entorno computacional, conviene observar cómo
estructuramos actividades cotidianas mediante tres patrones básicos:

\textbf{1. Flujo secuencial (cambio de un foco):}

\begin{enumerate}
\def\labelenumi{\arabic{enumi}.}
\tightlist
\item
  Comprobar si hay electricidad en la habitación.
\item
  Conseguir una escalera y un foco nuevo.
\item
  Apagar el interruptor de la luz.
\item
  Desenroscar el foco fundido girando hacia la izquierda.
\item
  Enroscar el foco nuevo girando hacia la derecha.
\item
  Encender el interruptor para verificar que funcione.
\item
  Concluir la tarea.
\end{enumerate}

En este caso, las acciones se suceden de forma estrictamente lineal,
ejecutando una instrucción detrás de otra.

\textbf{2. Toma de decisiones (preparación de café):}

\begin{enumerate}
\def\labelenumi{\arabic{enumi}.}
\item
  Servir agua caliente en una taza limpia.
\item
  Evaluar: ¿Desea azúcar?

  \begin{itemize}
  \tightlist
  \item
    \textbf{Sí:} Agregar una cucharada de azúcar y mezclar.
  \item
    \textbf{No:} Mezclar directamente.
  \end{itemize}
\item
  Concluir la preparación.
\end{enumerate}

En este flujo se introduce una bifurcación condicional basada en una
preferencia.

\textbf{3. Estructura iterativa (lavado de platos):}

\begin{enumerate}
\def\labelenumi{\arabic{enumi}.}
\item
  Tomar un plato sucio de la pila.
\item
  Enjabonar y enjuagar el plato.
\item
  Evaluar: ¿Quedan platos sucios en la pila?

  \begin{itemize}
  \tightlist
  \item
    \textbf{Sí:} Regresar al paso 1.
  \item
    \textbf{No:} Continuar al paso 4.
  \end{itemize}
\item
  Secarse las manos y finalizar.
\end{enumerate}

Al retornar al primer paso mientras existan platos pendientes, se
establece un ciclo de repetición.

\subsection{2.3 Variables, constantes, tipos de datos y
operadores}\label{variables-constantes-tipos-de-datos-y-operadores}

\textbf{Variables y constantes:}

\begin{itemize}
\tightlist
\item
  \textbf{Variables:} Espacios de memoria identificados con un nombre
  cuyo contenido puede cambiar durante la ejecución del programa (por
  ejemplo, \texttt{edad\ =\ 25} y más adelante \texttt{edad\ =\ 26}).
\item
  \textbf{Constantes:} Espacios de memoria cuyo valor se define una sola
  vez al inicio y permanece inalterable durante toda la ejecución (por
  ejemplo, \texttt{PI\ =\ 3.14159} o \texttt{TASA\_IVA\ =\ 0.16}).
\end{itemize}

\textbf{Buenas prácticas de nomenclatura:} Es recomendable emplear
nombres descriptivos que expliquen la función del identificador,
evitando letras sueltas y sin contexto. En la práctica profesional se
utilizan convenciones como \textbf{\texttt{snake\_case}}
(\texttt{calificacion\_final}) o \textbf{\texttt{camelCase}}
(\texttt{calificacionFinal}).

\textbf{Tipos de datos básicos:}

\begin{itemize}
\tightlist
\item
  \textbf{Numéricos:} Enteros (\texttt{15}) o decimales de punto
  flotante (\texttt{3.14}).
\item
  \textbf{Cadenas de caracteres (texto):} Delimitadas siempre entre
  comillas (\texttt{"Hola\ Mundo"}).
\end{itemize}

\textbf{Operadores:}

\begin{itemize}
\tightlist
\item
  \textbf{Aritméticos:} Suma (\texttt{+}), resta (\texttt{-}),
  multiplicación (\texttt{*}), división (\texttt{/}) y módulo o residuo
  (\texttt{MOD} o \texttt{\%}).
\item
  \textbf{Concatenación:} Permite unir texto y variables usando el signo
  \texttt{+} (por ejemplo, \texttt{"Tienes\ "\ +\ edad\ +\ "\ años"}).
\item
  \textbf{Lógicos:} Operadores \texttt{Y} (AND), \texttt{O} (OR) y
  \texttt{NO} (NOT) para articular condiciones compuestas.
\end{itemize}

\textbf{Diferencia clave: Asignación vs.~Comparación (\texttt{=} vs
\texttt{==}):}

\begin{itemize}
\tightlist
\item
  Un solo signo igual (\texttt{=}) se emplea para \textbf{asignar} un
  dato a una variable (por ejemplo, \texttt{saldo\ =\ 1000}).
\item
  El doble signo igual (\texttt{==}) se utiliza para \textbf{comparar}
  dos valores dentro de una condición lógica (por ejemplo,
  \texttt{SI\ opcion\ ==\ 1}).
\end{itemize}

\begin{center}\rule{0.5\linewidth}{0.5pt}\end{center}

\section{3. Construcción gradual de algoritmos y
DFDs}\label{construcciuxf3n-gradual-de-algoritmos-y-dfds}

Antes de trazar un diagrama, conviene seguir una metodología
estructurada: comprender el problema, identificar las entradas
requeridas, definir las operaciones o fórmulas necesarias y especificar
las salidas esperadas. En esta sección se aborda la construcción
progresiva de algoritmos mediante pseudocódigo, diagramas de flujo de
datos (DFD) y tablas de prueba de escritorio.

\textbf{Reglas para líneas de flujo y conectores:}

\begin{itemize}
\tightlist
\item
  El flujo natural se orienta de arriba hacia abajo y de izquierda a
  derecha.
\item
  Las líneas de flujo no deben cruzarse; cuando el diagrama se extiende
  o requiere enlazar puntos distantes, se utilizan conectores dentro de
  la misma página (círculo con identificador) o conectores fuera de
  página (pentágono).
\item
  A excepción del inicio y el fin, cada bloque debe tener al menos una
  línea de entrada y una de salida para evitar trayectorias inconclusas.
\end{itemize}

\textbf{La prueba de escritorio:} Consiste en una verificación manual
donde se simula la ejecución del algoritmo mediante una tabla,
registrando paso a paso el estado de las variables y las salidas
emitidas para comprobar su correcto funcionamiento y evaluar casos
extremos.

A continuación se ilustra la prueba de escritorio para un ciclo que
cuenta del 1 al 3 (\texttt{PARA\ i\ =\ 1\ HASTA\ 3}):

{\def\LTcaptype{none} % do not increment counter
\begin{longtable}[]{@{}
  >{\centering\arraybackslash}p{(\linewidth - 6\tabcolsep) * \real{0.2632}}
  >{\centering\arraybackslash}p{(\linewidth - 6\tabcolsep) * \real{0.2632}}
  >{\centering\arraybackslash}p{(\linewidth - 6\tabcolsep) * \real{0.2632}}
  >{\raggedright\arraybackslash}p{(\linewidth - 6\tabcolsep) * \real{0.2105}}@{}}
\toprule\noalign{}
\begin{minipage}[b]{\linewidth}\centering
Vuelta / Paso
\end{minipage} & \begin{minipage}[b]{\linewidth}\centering
¿Límite alcanzado?
\end{minipage} & \begin{minipage}[b]{\linewidth}\centering
Variable \texttt{i} (Memoria)
\end{minipage} & \begin{minipage}[b]{\linewidth}\raggedright
Acción en pantalla (Salida)
\end{minipage} \\
\midrule\noalign{}
\endhead
\bottomrule\noalign{}
\endlastfoot
Inicio & NO & 1 & ``Imprimiendo 1'' \\
Incremento & NO & 2 & ``Imprimiendo 2'' \\
Incremento & NO & 3 & ``Imprimiendo 3'' \\
Incremento & \textbf{SÍ (Límite)} & 4 & \emph{El programa sale del
ciclo} \\
\end{longtable}
}

\subsection{El esqueleto base (Estructura
vacía)}\label{el-esqueleto-base-estructura-vacuxeda}

Todo algoritmo requiere un punto de inicio y un punto de fin bien
delimitados. En los diagramas de flujo se representan mediante el
\textbf{óvalo} (\texttt{INICIO} y \texttt{FIN}), que formaliza el inicio
de la ejecución y la liberación final de los recursos del sistema.

\textbf{Pseudocódigo:}

\begin{verbatim}
INICIO
FIN
\end{verbatim}

\textbf{Diagrama de flujo:}

\begin{center}
\begin{adjustbox}{max width=\textwidth,max totalheight=.72\textheight,keepaspectratio}
\begin{tikzpicture}[flujo,node distance=5mm and 14mm]
  \node[iniciofin] (n0) {INICIO};
  \node[iniciofin,below=of n0] (n1) {FIN};
  \draw[draw=AzulLinea] (n0) -- (n1);
\end{tikzpicture}
\end{adjustbox}
\end{center}


\begin{center}\rule{0.5\linewidth}{0.5pt}\end{center}

\subsection{Algoritmos con instrucciones de
salida}\label{algoritmos-con-instrucciones-de-salida}

Para enviar información desde el programa hacia el usuario o hacia otros
dispositivos se utiliza el símbolo de salida, representado aquí con el
\textbf{trapecio de salida} \texttt{{[}\textbackslash{}\ ...\ /{]}}.
Esta instrucción despliega resultados en pantalla, genera impresiones o
envía mensajes de estado.

\subsubsection{Ejemplo 1: Saludo en
pantalla}\label{ejemplo-1-saludo-en-pantalla}

\textbf{Descripción:} Mostrar en pantalla un mensaje de bienvenida
estático con el texto \texttt{"Hola\ Alexander"}.

\textbf{Pseudocódigo:}

\begin{verbatim}
INICIO
    IMPRIMIR "Hola Alexander"
FIN
\end{verbatim}

\textbf{Diagrama de flujo:}

\begin{center}
\begin{adjustbox}{max width=\textwidth,max totalheight=.72\textheight,keepaspectratio}
\begin{tikzpicture}[flujo,node distance=5mm and 14mm]
  \node[iniciofin] (n0) {INICIO};
  \node[documento,below=of n0] (n1) {"Hola Alexander"};
  \node[iniciofin,below=of n1] (n2) {FIN};
  \draw[draw=AzulLinea] (n0) -- (n1);
  \draw[draw=TurquesaLinea] (n1) -- (n2);
\end{tikzpicture}
\end{adjustbox}
\end{center}


\textbf{Prueba de escritorio:}

{\def\LTcaptype{none} % do not increment counter
\begin{longtable}[]{@{}cl@{}}
\toprule\noalign{}
Paso & Acción en pantalla (Salida) \\
\midrule\noalign{}
\endhead
\bottomrule\noalign{}
\endlastfoot
Inicio & \emph{Ninguna} \\
Imprimir & ``Hola Alexander'' \\
Fin & \emph{Programa terminado} \\
\end{longtable}
}

\subsubsection{Ejemplo 2: Mensaje de estado de un
proceso}\label{ejemplo-2-mensaje-de-estado-de-un-proceso}

\textbf{Descripción:} Notificar al usuario que una tarea en segundo
plano ha comenzado mediante la impresión de un mensaje informativo.

\textbf{Pseudocódigo:}

\begin{verbatim}
INICIO
    IMPRIMIR "Iniciando entrenamiento del modelo..."
FIN
\end{verbatim}

\textbf{Diagrama de flujo:}

\begin{center}
\begin{adjustbox}{max width=\textwidth,max totalheight=.72\textheight,keepaspectratio}
\begin{tikzpicture}[flujo,node distance=5mm and 14mm]
  \node[iniciofin] (n0) {INICIO};
  \node[documento,below=of n0,text width=45mm] (n1) {"Iniciando entrenamiento del modelo..."};
  \node[iniciofin,below=of n1] (n2) {FIN};
  \draw[draw=AzulLinea] (n0) -- (n1);
  \draw[draw=TurquesaLinea] (n1) -- (n2);
\end{tikzpicture}
\end{adjustbox}
\end{center}


\textbf{Prueba de escritorio:}

{\def\LTcaptype{none} % do not increment counter
\begin{longtable}[]{@{}cl@{}}
\toprule\noalign{}
Paso & Acción en pantalla (Salida) \\
\midrule\noalign{}
\endhead
\bottomrule\noalign{}
\endlastfoot
Inicio & \emph{Ninguna} \\
Imprimir & ``Iniciando entrenamiento del modelo\ldots{}'' \\
Fin & \emph{Programa terminado} \\
\end{longtable}
}

\subsubsection{Ejemplo 3: Emisión de un código de
error}\label{ejemplo-3-emisiuxf3n-de-un-cuxf3digo-de-error}

\textbf{Descripción:} Mostrar un código de error estandarizado en
pantalla cuando no se localiza un recurso o archivo necesario.

\textbf{Pseudocódigo:}

\begin{verbatim}
INICIO
    IMPRIMIR "Error 404: Dataset no encontrado"
FIN
\end{verbatim}

\textbf{Diagrama de flujo:}

\begin{center}
\begin{adjustbox}{max width=\textwidth,max totalheight=.72\textheight,keepaspectratio}
\begin{tikzpicture}[flujo,node distance=5mm and 14mm]
  \node[iniciofin] (n0) {INICIO};
  \node[documento,below=of n0] (n1) {"Error 404: Dataset no encontrado"};
  \node[iniciofin,below=of n1] (n2) {FIN};
  \draw[draw=AzulLinea] (n0) -- (n1);
  \draw[draw=TurquesaLinea] (n1) -- (n2);
\end{tikzpicture}
\end{adjustbox}
\end{center}


\textbf{Prueba de escritorio:}

{\def\LTcaptype{none} % do not increment counter
\begin{longtable}[]{@{}cl@{}}
\toprule\noalign{}
Paso & Acción en pantalla (Salida) \\
\midrule\noalign{}
\endhead
\bottomrule\noalign{}
\endlastfoot
Inicio & \emph{Ninguna} \\
Imprimir & ``Error 404: Dataset no encontrado'' \\
Fin & \emph{Programa terminado} \\
\end{longtable}
}

\begin{center}\rule{0.5\linewidth}{0.5pt}\end{center}

\subsection{Algoritmos con captura de datos (Entrada y
Salida)}\label{algoritmos-con-captura-de-datos-entrada-y-salida}

Para recibir información proporcionada por el usuario o por sensores se
utiliza el \textbf{paralelogramo} \texttt{{[}\ /\ ...\ /\ {]}}, que
representa la operación de lectura (\texttt{LEER}) y almacena los datos
ingresados en variables de memoria.

\subsubsection{Ejemplo 1: Lectura y confirmación de
edad}\label{ejemplo-1-lectura-y-confirmaciuxf3n-de-edad}

\textbf{Descripción:} Solicitar la edad de un paciente y mostrarla en
pantalla junto con un mensaje de confirmación.

\textbf{Pseudocódigo:}

\begin{verbatim}
INICIO
    LEER edad
    IMPRIMIR "Tu edad registrada es: " + edad
FIN
\end{verbatim}

\textbf{Diagrama de flujo:}

\begin{center}
\begin{adjustbox}{max width=\textwidth,max totalheight=.72\textheight,keepaspectratio}
\begin{tikzpicture}[flujo,node distance=5mm and 14mm]
  \node[iniciofin] (n0) {INICIO};
  \node[entradaSalida,below=of n0] (n1) {edad};
  \node[documento,below=of n1] (n2) {"Tu edad registrada es: " + edad};
  \node[iniciofin,below=of n2] (n3) {FIN};
  \draw[draw=AzulLinea] (n0) -- (n1);
  \draw[draw=TurquesaLinea] (n1) -- (n2);
  \draw[draw=RosaLinea] (n2) -- (n3);
\end{tikzpicture}
\end{adjustbox}
\end{center}


\textbf{Prueba de escritorio:}

{\def\LTcaptype{none} % do not increment counter
\begin{longtable}[]{@{}ccl@{}}
\toprule\noalign{}
Paso & Variable \texttt{edad} & Acción en pantalla (Salida) \\
\midrule\noalign{}
\endhead
\bottomrule\noalign{}
\endlastfoot
Leer & \textbf{25} \emph{(ingreso)} & \emph{Esperando entrada del
teclado} \\
Imprimir & 25 & ``Tu edad registrada es: 25'' \\
\end{longtable}
}

\subsubsection{Ejemplo 2: Concatenación de fecha y
ciudad}\label{ejemplo-2-concatenaciuxf3n-de-fecha-y-ciudad}

\textbf{Descripción:} Solicitar el nombre de una ciudad y el año en
curso para construir un encabezado formal mediante concatenación de
texto.

\textbf{Pseudocódigo:}

\begin{verbatim}
INICIO
    LEER ciudad
    LEER anio
    IMPRIMIR ciudad + ", Col. a 18 de agosto de " + anio
FIN
\end{verbatim}

\textbf{Diagrama de flujo:}

\begin{center}
\begin{adjustbox}{max width=\textwidth,max totalheight=.72\textheight,keepaspectratio}
\begin{tikzpicture}[flujo,node distance=5mm and 14mm]
  \node[iniciofin] (n0) {INICIO};
  \node[entradaSalida,below=of n0] (n1) {ciudad};
  \node[entradaSalida,below=of n1] (n2) {año};
  \node[documento,below=of n2,text width=45mm] (n3) {ciudad + ", Col. a 18 de agosto de " + año};
  \node[iniciofin,below=of n3] (n4) {FIN};
  \draw[draw=AzulLinea] (n0) -- (n1);
  \draw[draw=TurquesaLinea] (n1) -- (n2);
  \draw[draw=RosaLinea] (n2) -- (n3);
  \draw[draw=MoradoLinea] (n3) -- (n4);
\end{tikzpicture}
\end{adjustbox}
\end{center}


\textbf{Prueba de escritorio:}

{\def\LTcaptype{none} % do not increment counter
\begin{longtable}[]{@{}cccl@{}}
\toprule\noalign{}
Paso & Var \texttt{ciudad} & Var \texttt{anio} & Acción en pantalla
(Salida) \\
\midrule\noalign{}
\endhead
\bottomrule\noalign{}
\endlastfoot
Leer & \textbf{``Colima''} & - & \emph{Esperando texto} \\
Leer & ``Colima'' & \textbf{2026} & \emph{Esperando número} \\
Imprimir & ``Colima'' & 2026 & ``Colima, Col. a 18 de agosto de
2026'' \\
\end{longtable}
}

\subsubsection{Ejemplo 3: Carga de modelo por
nombre}\label{ejemplo-3-carga-de-modelo-por-nombre}

\textbf{Descripción:} Solicitar el nombre de un modelo predictivo e
imprimir un mensaje confirmando que se encuentra listo para su
utilización.

\textbf{Pseudocódigo:}

\begin{verbatim}
INICIO
    LEER nombre_modelo
    IMPRIMIR "El modelo " + nombre_modelo + " está listo para usarse."
FIN
\end{verbatim}

\textbf{Diagrama de flujo:}

\begin{center}
\begin{adjustbox}{max width=\textwidth,max totalheight=.72\textheight,keepaspectratio}
\begin{tikzpicture}[flujo,node distance=5mm and 14mm]
  \node[iniciofin] (n0) {INICIO};
  \node[entradaSalida,below=of n0] (n1) {nombre\_modelo};
  \node[documento,below=of n1,text width=45mm] (n2) {"El modelo " + nombre\_modelo + " está listo para usarse"};
  \node[iniciofin,below=of n2] (n3) {FIN};
  \draw[draw=AzulLinea] (n0) -- (n1);
  \draw[draw=TurquesaLinea] (n1) -- (n2);
  \draw[draw=RosaLinea] (n2) -- (n3);
\end{tikzpicture}
\end{adjustbox}
\end{center}


\textbf{Prueba de escritorio:}

{\def\LTcaptype{none} % do not increment counter
\begin{longtable}[]{@{}
  >{\centering\arraybackslash}p{(\linewidth - 4\tabcolsep) * \real{0.3571}}
  >{\centering\arraybackslash}p{(\linewidth - 4\tabcolsep) * \real{0.3571}}
  >{\raggedright\arraybackslash}p{(\linewidth - 4\tabcolsep) * \real{0.2857}}@{}}
\toprule\noalign{}
\begin{minipage}[b]{\linewidth}\centering
Paso
\end{minipage} & \begin{minipage}[b]{\linewidth}\centering
Var \texttt{nombre\_modelo}
\end{minipage} & \begin{minipage}[b]{\linewidth}\raggedright
Acción en pantalla (Salida)
\end{minipage} \\
\midrule\noalign{}
\endhead
\bottomrule\noalign{}
\endlastfoot
Leer & \textbf{``Regresión''} & \emph{Esperando texto} \\
Imprimir & ``Regresión'' & ``El modelo Regresión está listo para
usarse.'' \\
\end{longtable}
}

\begin{center}\rule{0.5\linewidth}{0.5pt}\end{center}

\subsection{Algoritmos con Entrada, Proceso y
Salida}\label{algoritmos-con-entrada-proceso-y-salida}

El \textbf{rectángulo} \texttt{{[}\ {]}} representa operaciones de
proceso, tales como cálculos aritméticos, asignaciones y
transformaciones de datos que modifican el estado de las variables en
memoria.

\subsubsection{Ejemplo 1: Cálculo de edad según el año de
nacimiento}\label{ejemplo-1-cuxe1lculo-de-edad-seguxfan-el-auxf1o-de-nacimiento}

\textbf{Descripción:} Solicitar el año de nacimiento del usuario, restar
dicho valor del año actual (2026) y mostrar la edad resultante.

\textbf{Pseudocódigo:}

\begin{verbatim}
INICIO
    LEER anio_nacimiento
    edad_calculada = 2026 - anio_nacimiento
    IMPRIMIR "Tienes " + edad_calculada + " años."
FIN
\end{verbatim}

\textbf{Diagrama de flujo:}

\begin{center}
\begin{adjustbox}{max width=\textwidth,max totalheight=.72\textheight,keepaspectratio}
\begin{tikzpicture}[flujo,node distance=5mm and 14mm]
  \node[iniciofin] (n0) {INICIO};
  \node[entradaSalida,below=of n0] (n1) {anio\_nacimiento};
  \node[proceso,below=of n1,text width=45mm] (n2) {edad\_calculada = 2026 - anio\_nacimiento};
  \node[documento,below=of n2,text width=45mm] (n3) {"Tienes " + edad\_calculada + " años"};
  \node[iniciofin,below=of n3] (n4) {FIN};
  \draw[draw=AzulLinea] (n0) -- (n1);
  \draw[draw=TurquesaLinea] (n1) -- (n2);
  \draw[draw=RosaLinea] (n2) -- (n3);
  \draw[draw=MoradoLinea] (n3) -- (n4);
\end{tikzpicture}
\end{adjustbox}
\end{center}


\textbf{Prueba de escritorio:}

{\def\LTcaptype{none} % do not increment counter
\begin{longtable}[]{@{}
  >{\centering\arraybackslash}p{(\linewidth - 6\tabcolsep) * \real{0.2632}}
  >{\centering\arraybackslash}p{(\linewidth - 6\tabcolsep) * \real{0.2632}}
  >{\centering\arraybackslash}p{(\linewidth - 6\tabcolsep) * \real{0.2632}}
  >{\raggedright\arraybackslash}p{(\linewidth - 6\tabcolsep) * \real{0.2105}}@{}}
\toprule\noalign{}
\begin{minipage}[b]{\linewidth}\centering
Paso
\end{minipage} & \begin{minipage}[b]{\linewidth}\centering
Var \texttt{anio\_nacimiento}
\end{minipage} & \begin{minipage}[b]{\linewidth}\centering
Var \texttt{edad\_calculada}
\end{minipage} & \begin{minipage}[b]{\linewidth}\raggedright
Acción en pantalla (Salida)
\end{minipage} \\
\midrule\noalign{}
\endhead
\bottomrule\noalign{}
\endlastfoot
Leer & \textbf{2000} & - & \emph{Esperando año numérico} \\
Proceso & 2000 & \textbf{26} \emph{(2026-2000)} & \emph{Operación
interna en memoria} \\
Imprimir & 2000 & 26 & ``Tienes 26 años.'' \\
\end{longtable}
}

\subsubsection{Ejemplo 2: Promedio de tres
calificaciones}\label{ejemplo-2-promedio-de-tres-calificaciones}

\textbf{Descripción:} Leer tres notas parciales, calcular su promedio
aritmético mediante suma y división, y mostrar el promedio final.

\textbf{Pseudocódigo:}

\begin{verbatim}
INICIO
    LEER calif_1
    LEER calif_2
    LEER calif_3
    promedio = (calif_1 + calif_2 + calif_3) / 3
    IMPRIMIR "El promedio final es: " + promedio
FIN
\end{verbatim}

\textbf{Diagrama de flujo:}

\begin{center}
\begin{adjustbox}{max width=\textwidth,max totalheight=.72\textheight,keepaspectratio}
\begin{tikzpicture}[flujo,node distance=5mm and 14mm]
  \node[iniciofin] (n0) {INICIO};
  \node[entradaSalida,below=of n0] (n1) {calif\_1, calif\_2, calif\_3};
  \node[proceso,below=of n1,text width=45mm] (n2) {promedio = (calif\_1 + calif\_2 + calif\_3) / 3};
  \node[documento,below=of n2] (n3) {"El promedio final es: " + promedio};
  \node[iniciofin,below=of n3] (n4) {FIN};
  \draw[draw=AzulLinea] (n0) -- (n1);
  \draw[draw=TurquesaLinea] (n1) -- (n2);
  \draw[draw=RosaLinea] (n2) -- (n3);
  \draw[draw=MoradoLinea] (n3) -- (n4);
\end{tikzpicture}
\end{adjustbox}
\end{center}


\textbf{Prueba de escritorio:}

{\def\LTcaptype{none} % do not increment counter
\begin{longtable}[]{@{}cccccl@{}}
\toprule\noalign{}
Paso & \texttt{calif\_1} & \texttt{calif\_2} & \texttt{calif\_3} &
\texttt{promedio} & Pantalla \\
\midrule\noalign{}
\endhead
\bottomrule\noalign{}
\endlastfoot
Leer & \textbf{8} & \textbf{9} & \textbf{10} & - & \emph{Esperando 3
datos} \\
Proceso & 8 & 9 & 10 & \textbf{9} & \emph{Cálculo interno} \\
Imprimir & 8 & 9 & 10 & 9 & ``El promedio final es: 9'' \\
\end{longtable}
}

\subsubsection{Ejemplo 3: Área de un terreno
rectangular}\label{ejemplo-3-uxe1rea-de-un-terreno-rectangular}

\textbf{Descripción:} Solicitar la base y la altura de un terreno
rectangular, multiplicar ambas dimensiones y mostrar el área en metros
cuadrados.

\textbf{Pseudocódigo:}

\begin{verbatim}
INICIO
    LEER base
    LEER altura
    area = base * altura
    IMPRIMIR "El área total es de: " + area + " metros cuadrados"
FIN
\end{verbatim}

\textbf{Diagrama de flujo:}

\begin{center}
\begin{adjustbox}{max width=\textwidth,max totalheight=.72\textheight,keepaspectratio}
\begin{tikzpicture}[flujo,node distance=5mm and 14mm]
  \node[iniciofin] (n0) {INICIO};
  \node[entradaSalida,below=of n0] (n1) {base, altura};
  \node[proceso,below=of n1] (n2) {area = base * altura};
  \node[documento,below=of n2,text width=45mm] (n3) {"El área total es de: " + area + " metros cuadrados"};
  \node[iniciofin,below=of n3] (n4) {FIN};
  \draw[draw=AzulLinea] (n0) -- (n1);
  \draw[draw=TurquesaLinea] (n1) -- (n2);
  \draw[draw=RosaLinea] (n2) -- (n3);
  \draw[draw=MoradoLinea] (n3) -- (n4);
\end{tikzpicture}
\end{adjustbox}
\end{center}


\textbf{Prueba de escritorio:}

{\def\LTcaptype{none} % do not increment counter
\begin{longtable}[]{@{}ccccl@{}}
\toprule\noalign{}
Paso & Var \texttt{base} & Var \texttt{altura} & Var \texttt{area} &
Acción en pantalla \\
\midrule\noalign{}
\endhead
\bottomrule\noalign{}
\endlastfoot
Leer & \textbf{10} & \textbf{5} & - & \emph{Esperando 2 datos} \\
Proceso & 10 & 5 & \textbf{50} & \emph{Cálculo interno} \\
Imprimir & 10 & 5 & 50 & ``El área total es de: 50 metros cuadrados'' \\
\end{longtable}
}

\begin{center}\rule{0.5\linewidth}{0.5pt}\end{center}

\section{4. Estructuras condicionales (toma de
decisiones)}\label{estructuras-condicionales-toma-de-decisiones}

El \textbf{rombo} \texttt{\{\ \}} evalúa una expresión lógica que solo
puede resultar verdadera o falsa, bifurcando el flujo del programa en
dos caminos independientes según el resultado. A continuación se
presentan casos con decisiones simples, anidadas y combinadas mediante
operadores lógicos.

\subsection{Condicional simple: Par o
Impar}\label{condicional-simple-par-o-impar}

\textbf{Descripción:} Determinar si un número entero es par o impar
evaluando si el residuo de su división entre 2 (\texttt{MOD}) es igual a
cero.

\textbf{Pseudocódigo:}

\begin{verbatim}
INICIO
    LEER numero
    residuo = numero MOD 2
    SI residuo == 0 ENTONCES
        IMPRIMIR "El número es Par"
    SINO
        IMPRIMIR "El número es Impar"
    FIN SI
FIN
\end{verbatim}

\begin{center}
\begin{adjustbox}{max width=\textwidth,max totalheight=.72\textheight,keepaspectratio}
\begin{tikzpicture}[flujo,node distance=5mm and 14mm]
  \node[iniciofin] (n0) {INICIO};
  \node[entradaSalida,below=of n0] (n1) {número};
  \node[proceso,below=of n1] (n2) {residuo = número MOD 2};
  \node[decision,below=of n2] (n3) {residuo == 0?};
  \node[documento,below=of n3] (n4) {"El número es Par"};
  \node[documento,below=of n4] (n5) {"El número es Impar"};
  \node[iniciofin,below=of n5] (n6) {FIN};
  \draw[draw=AzulLinea] (n0) -- (n1);
  \draw[draw=TurquesaLinea] (n1) -- (n2);
  \draw[draw=RosaLinea] (n2) -- (n3);
  \draw[draw=VerdeLinea] (n3) -- node[etiqueta,right=2pt,pos=.45] {YES} (n4);
  \draw[draw=NaranjaLinea] (n3.east) --  node[etiqueta,above=2pt,pos=.45] {NO} ++(43mm,0) |-  (n5.east);
  \draw[draw=TurquesaLinea] (n4.east) --  ++(48mm,0) |-  (n6.east);
  \draw[draw=RosaLinea] (n5) -- (n6);
\end{tikzpicture}
\end{adjustbox}
\end{center}


\textbf{Prueba de escritorio (evaluando el valor impar 7):}

{\def\LTcaptype{none} % do not increment counter
\begin{longtable}[]{@{}ccccl@{}}
\toprule\noalign{}
Paso & \texttt{numero} & \texttt{residuo} & Condición
(\texttt{residuo\ ==\ 0}) & Pantalla \\
\midrule\noalign{}
\endhead
\bottomrule\noalign{}
\endlastfoot
Leer & \textbf{7} & - & - & \emph{Esperando dato} \\
Proceso & 7 & \textbf{1} & - & \emph{Cálculo interno} \\
Rombo & 7 & 1 & \textbf{NO (Falso)} & \emph{Bifurca al Sino} \\
Imprimir & 7 & 1 & - & ``El número es Impar'' \\
\end{longtable}
}

\subsection{Condicionales anidados: Clasificar un
número}\label{condicionales-anidados-clasificar-un-nuxfamero}

\textbf{Descripción:} Clasificar un número ingresado indicando si es
positivo, negativo o exactamente cero mediante dos decisiones
consecutivas.

\textbf{Pseudocódigo:}

\begin{verbatim}
INICIO
    LEER valor
    SI valor == 0 ENTONCES
        IMPRIMIR "El valor es Cero"
    SINO
        SI valor > 0 ENTONCES
            IMPRIMIR "Es Positivo"
        SINO
            IMPRIMIR "Es Negativo"
        FIN SI
    FIN SI
FIN
\end{verbatim}

\begin{center}
\begin{adjustbox}{max width=\textwidth,max totalheight=.72\textheight,keepaspectratio}
\begin{tikzpicture}[flujo,node distance=5mm and 14mm]
  \node[iniciofin] (n0) {INICIO};
  \node[entradaSalida,below=of n0] (n1) {valor};
  \node[decision,below=of n1] (n2) {valor == 0?};
  \node[documento,below=of n2] (n3) {"El valor es Cero"};
  \node[decision,below=of n3] (n4) {valor > 0?};
  \node[documento,below=of n4] (n5) {"Es Positivo"};
  \node[documento,below=of n5] (n6) {"Es Negativo"};
  \node[iniciofin,below=of n6] (n7) {FIN};
  \draw[draw=AzulLinea] (n0) -- (n1);
  \draw[draw=TurquesaLinea] (n1) -- (n2);
  \draw[draw=VerdeLinea] (n2) -- node[etiqueta,right=2pt,pos=.45] {YES} (n3);
  \draw[draw=NaranjaLinea] (n2.east) --  node[etiqueta,above=2pt,pos=.45] {NO} ++(43mm,0) |-  (n4.east);
  \draw[draw=VerdeLinea] (n4) -- node[etiqueta,right=2pt,pos=.45] {YES} (n5);
  \draw[draw=NaranjaLinea] (n4.east) --  node[etiqueta,above=2pt,pos=.45] {NO} ++(48mm,0) |-  (n6.east);
  \draw[draw=RosaLinea] (n3.east) --  ++(38mm,0) |-  (n7.east);
  \draw[draw=MoradoLinea] (n5.east) --  ++(43mm,0) |-  (n7.east);
  \draw[draw=AzulLinea] (n6) -- (n7);
\end{tikzpicture}
\end{adjustbox}
\end{center}


\textbf{Prueba de escritorio (evaluando el valor negativo -5):}

{\def\LTcaptype{none} % do not increment counter
\begin{longtable}[]{@{}
  >{\centering\arraybackslash}p{(\linewidth - 8\tabcolsep) * \real{0.2083}}
  >{\centering\arraybackslash}p{(\linewidth - 8\tabcolsep) * \real{0.2083}}
  >{\centering\arraybackslash}p{(\linewidth - 8\tabcolsep) * \real{0.2083}}
  >{\centering\arraybackslash}p{(\linewidth - 8\tabcolsep) * \real{0.2083}}
  >{\raggedright\arraybackslash}p{(\linewidth - 8\tabcolsep) * \real{0.1667}}@{}}
\toprule\noalign{}
\begin{minipage}[b]{\linewidth}\centering
Paso
\end{minipage} & \begin{minipage}[b]{\linewidth}\centering
\texttt{valor}
\end{minipage} & \begin{minipage}[b]{\linewidth}\centering
Condición 1 (\texttt{valor\ ==\ 0})
\end{minipage} & \begin{minipage}[b]{\linewidth}\centering
Condición 2 (\texttt{valor\ \textgreater{}\ 0})
\end{minipage} & \begin{minipage}[b]{\linewidth}\raggedright
Pantalla
\end{minipage} \\
\midrule\noalign{}
\endhead
\bottomrule\noalign{}
\endlastfoot
Leer & \textbf{-5} & - & - & \emph{Esperando dato} \\
Rombo 1 & -5 & \textbf{NO (Falso)} & - & \emph{Bifurca a la rama
alternativa} \\
Rombo 2 & -5 & - & \textbf{NO (Falso)} & \emph{Bifurca al Sino
interno} \\
Imprimir & -5 & - & - & ``Es Negativo'' \\
\end{longtable}
}

\subsection{Operadores lógicos (Y, O, NO): Evaluación
conjunta}\label{operadores-luxf3gicos-y-o-no-evaluaciuxf3n-conjunta}

\textbf{Descripción:} Aprobar a un estudiante únicamente cuando su
calificación sea mayor o igual a 6 y su asistencia supere el 80\%,
combinando ambas condiciones con el operador \texttt{Y}.

\textbf{Pseudocódigo:}

\begin{verbatim}
INICIO
    LEER calificacion
    LEER asistencia
    SI calificacion >= 6 Y asistencia > 80 ENTONCES
        IMPRIMIR "Aprobado"
    SINO
        IMPRIMIR "Reprobado"
    FIN SI
FIN
\end{verbatim}

\begin{center}
\begin{adjustbox}{max width=\textwidth,max totalheight=.72\textheight,keepaspectratio}
\begin{tikzpicture}[flujo,node distance=5mm and 14mm]
  \node[iniciofin] (n0) {INICIO};
  \node[entradaSalida,below=of n0] (n1) {calif, asistencia};
  \node[decision,below=of n1,text width=45mm] (n2) {calif >= 6\\Y\\asistencia > 80?};
  \node[documento,below=of n2] (n3) {"Aprobado"};
  \node[documento,below=of n3] (n4) {"Reprobado"};
  \node[iniciofin,below=of n4] (n5) {FIN};
  \draw[draw=AzulLinea] (n0) -- (n1);
  \draw[draw=TurquesaLinea] (n1) -- (n2);
  \draw[draw=VerdeLinea] (n2) -- node[etiqueta,right=2pt,pos=.45] {YES} (n3);
  \draw[draw=NaranjaLinea] (n2.east) --  node[etiqueta,above=2pt,pos=.45] {NO} ++(43mm,0) |-  (n4.east);
  \draw[draw=AzulLinea] (n3.east) --  ++(48mm,0) |-  (n5.east);
  \draw[draw=TurquesaLinea] (n4) -- (n5);
\end{tikzpicture}
\end{adjustbox}
\end{center}


\textbf{Prueba de escritorio (calificación aprobatoria pero asistencia
insuficiente: 9 y 70):}

{\def\LTcaptype{none} % do not increment counter
\begin{longtable}[]{@{}
  >{\centering\arraybackslash}p{(\linewidth - 8\tabcolsep) * \real{0.2083}}
  >{\centering\arraybackslash}p{(\linewidth - 8\tabcolsep) * \real{0.2083}}
  >{\centering\arraybackslash}p{(\linewidth - 8\tabcolsep) * \real{0.2083}}
  >{\centering\arraybackslash}p{(\linewidth - 8\tabcolsep) * \real{0.2083}}
  >{\raggedright\arraybackslash}p{(\linewidth - 8\tabcolsep) * \real{0.1667}}@{}}
\toprule\noalign{}
\begin{minipage}[b]{\linewidth}\centering
Paso
\end{minipage} & \begin{minipage}[b]{\linewidth}\centering
\texttt{calif}
\end{minipage} & \begin{minipage}[b]{\linewidth}\centering
\texttt{asistencia}
\end{minipage} & \begin{minipage}[b]{\linewidth}\centering
Condición lógica conjunta (\texttt{Y})
\end{minipage} & \begin{minipage}[b]{\linewidth}\raggedright
Pantalla
\end{minipage} \\
\midrule\noalign{}
\endhead
\bottomrule\noalign{}
\endlastfoot
Leer & \textbf{9} & \textbf{70} & - & \emph{Esperando datos} \\
Rombo & 9 & 70 & \textbf{NO (Falso, pues 70 no es \textgreater{} 80)} &
\emph{Bifurca al Sino} \\
Imprimir & 9 & 70 & - & ``Reprobado'' \\
\end{longtable}
}

\begin{center}\rule{0.5\linewidth}{0.5pt}\end{center}

\section{5. Estructuras iterativas
(ciclos)}\label{estructuras-iterativas-ciclos}

Las estructuras iterativas permiten repetir un bloque de instrucciones
de manera controlada. Según el momento en que se evalúa la condición de
permanencia y el conocimiento previo del número de repeticiones, se
dividen en tres tipos principales.

\subsection{5.1 El ciclo MIENTRAS (bucle
While)}\label{el-ciclo-mientras-bucle-while}

Evalúa la condición al \textbf{inicio}, antes de ejecutar las
instrucciones internas: - \textbf{Comportamiento:} Si la condición es
verdadera, ejecuta el cuerpo del ciclo; si es falsa desde el primer
instante, no entra al bloque. - \textbf{Caso de uso:} Procesos donde el
número total de repeticiones no se conoce de antemano y depende de una
condición dinámica.

\subsubsection{Ejercicio MIENTRAS 1: Promedio de N
números}\label{ejercicio-mientras-1-promedio-de-n-nuxfameros}

\textbf{Descripción:} Solicitar al usuario la cantidad total de números
que desea ingresar, acumular sus valores mediante un ciclo y calcular el
promedio final.

\textbf{Pseudocódigo:}

\begin{verbatim}
INICIO
    LEER cantidad_n
    contador = 1
    suma = 0
    MIENTRAS contador <= cantidad_n HACER
        LEER numero_actual
        suma = suma + numero_actual
        contador = contador + 1
    FIN MIENTRAS
    promedio = suma / cantidad_n
    IMPRIMIR "La media es: " + promedio
FIN
\end{verbatim}

\begin{center}
\begin{adjustbox}{max width=\textwidth,max totalheight=.72\textheight,keepaspectratio}
\begin{tikzpicture}[flujo,node distance=5mm and 14mm]
  \node[iniciofin] (n0) {INICIO};
  \node[entradaSalida,below=of n0] (n1) {cantidad\_n};
  \node[proceso,below=of n1,text width=45mm] (n2) {contador = 1\\suma = 0};
  \node[decision,below=of n2] (n3) {contador <= cantidad\_n?};
  \node[entradaSalida,below=of n3] (n4) {numero\_actual};
  \node[proceso,below=of n4,text width=45mm] (n5) {suma = suma + numero\_actual\\contador = contador + 1};
  \node[proceso,below=of n5] (n6) {promedio = suma / cantidad\_n};
  \node[documento,below=of n6] (n7) {"La media es: " + promedio};
  \node[iniciofin,below=of n7] (n8) {FIN};
  \draw[draw=AzulLinea] (n0) -- (n1);
  \draw[draw=TurquesaLinea] (n1) -- (n2);
  \draw[draw=RosaLinea] (n2) -- (n3);
  \draw[draw=VerdeLinea] (n3) -- node[etiqueta,right=2pt,pos=.45] {YES} (n4);
  \draw[draw=AzulLinea] (n4) -- (n5);
  \draw[draw=MoradoLinea] (n5.west) --  ++(-43mm,0) |- ($(n2.south)!.5!(n3.north)$);
  \draw[draw=NaranjaLinea] (n3.east) --  node[etiqueta,above=2pt,pos=.45] {NO} ++(43mm,0) |-  (n6.east);
  \draw[draw=MoradoLinea] (n6) -- (n7);
  \draw[draw=AzulLinea] (n7) -- (n8);
\end{tikzpicture}
\end{adjustbox}
\end{center}


\textbf{Prueba de escritorio (promediando 2 números: 8 y 10):}

{\def\LTcaptype{none} % do not increment counter
\begin{longtable}[]{@{}
  >{\centering\arraybackslash}p{(\linewidth - 12\tabcolsep) * \real{0.1471}}
  >{\centering\arraybackslash}p{(\linewidth - 12\tabcolsep) * \real{0.1471}}
  >{\centering\arraybackslash}p{(\linewidth - 12\tabcolsep) * \real{0.1471}}
  >{\centering\arraybackslash}p{(\linewidth - 12\tabcolsep) * \real{0.1471}}
  >{\centering\arraybackslash}p{(\linewidth - 12\tabcolsep) * \real{0.1471}}
  >{\centering\arraybackslash}p{(\linewidth - 12\tabcolsep) * \real{0.1471}}
  >{\raggedright\arraybackslash}p{(\linewidth - 12\tabcolsep) * \real{0.1176}}@{}}
\toprule\noalign{}
\begin{minipage}[b]{\linewidth}\centering
Vuelta
\end{minipage} & \begin{minipage}[b]{\linewidth}\centering
\texttt{cantidad\_n}
\end{minipage} & \begin{minipage}[b]{\linewidth}\centering
\texttt{contador}
\end{minipage} & \begin{minipage}[b]{\linewidth}\centering
\texttt{suma}
\end{minipage} & \begin{minipage}[b]{\linewidth}\centering
\texttt{numero\_actual}
\end{minipage} & \begin{minipage}[b]{\linewidth}\centering
Condición (\texttt{\textless{}=\ n})
\end{minipage} & \begin{minipage}[b]{\linewidth}\raggedright
Pantalla
\end{minipage} \\
\midrule\noalign{}
\endhead
\bottomrule\noalign{}
\endlastfoot
Inicio & 2 & 1 & 0 & - & - & \emph{Esperando N} \\
1 & 2 & 1 & 0 & \textbf{8} & \textbf{SÍ (1 \textless= 2)} &
\emph{Esperando número} \\
1 (Fin) & 2 & \textbf{2} & \textbf{8} & 8 & - & - \\
2 & 2 & 2 & 8 & \textbf{10} & \textbf{SÍ (2 \textless= 2)} &
\emph{Esperando número} \\
2 (Fin) & 2 & \textbf{3} & \textbf{18} & 10 & - & - \\
3 & 2 & 3 & 18 & - & \textbf{NO (3 \textless= 2 es Falso)} & \emph{Sale
del ciclo} \\
Salida & 2 & 3 & 18 & - & - & ``La media es: 9'' \\
\end{longtable}
}

\subsubsection{Ejercicio MIENTRAS 2: Conteo e impresión de pares del 2
al
10}\label{ejercicio-mientras-2-conteo-e-impresiuxf3n-de-pares-del-2-al-10}

\textbf{Descripción:} Desplegar los números pares entre 2 y 10
incrementando de dos en dos, y mostrar al finalizar el total de números
impresos.

\textbf{Pseudocódigo:}

\begin{verbatim}
INICIO
    numero = 2
    cantidad_pares = 0
    MIENTRAS numero <= 10 HACER
        IMPRIMIR "Encontré el par: " + numero
        cantidad_pares = cantidad_pares + 1
        numero = numero + 2
    FIN MIENTRAS
    IMPRIMIR "Total de pares: " + cantidad_pares
FIN
\end{verbatim}

\begin{center}
\begin{adjustbox}{max width=\textwidth,max totalheight=.72\textheight,keepaspectratio}
\begin{tikzpicture}[flujo,node distance=5mm and 14mm]
  \node[iniciofin] (n0) {INICIO};
  \node[proceso,below=of n0,text width=45mm] (n1) {número = 2\\cantidad\_pares = 0};
  \node[decision,below=of n1] (n2) {número <= 10?};
  \node[documento,below=of n2] (n3) {"Encontré el par: " + número};
  \node[proceso,below=of n3,text width=45mm] (n4) {cantidad\_pares = cantidad\_pares + 1\\número = número + 2};
  \node[documento,below=of n4] (n5) {"Total de pares: " + cantidad\_pares};
  \node[iniciofin,below=of n5] (n6) {FIN};
  \draw[draw=AzulLinea] (n0) -- (n1);
  \draw[draw=TurquesaLinea] (n1) -- (n2);
  \draw[draw=VerdeLinea] (n2) -- node[etiqueta,right=2pt,pos=.45] {YES} (n3);
  \draw[draw=MoradoLinea] (n3) -- (n4);
  \draw[draw=MoradoLinea] (n4.west) --  ++(-43mm,0) |- ($(n1.south)!.5!(n2.north)$);
  \draw[draw=NaranjaLinea] (n2.east) --  node[etiqueta,above=2pt,pos=.45] {NO} ++(43mm,0) |-  (n5.east);
  \draw[draw=RosaLinea] (n5) -- (n6);
\end{tikzpicture}
\end{adjustbox}
\end{center}


\textbf{Prueba de escritorio:}

{\def\LTcaptype{none} % do not increment counter
\begin{longtable}[]{@{}
  >{\centering\arraybackslash}p{(\linewidth - 8\tabcolsep) * \real{0.2083}}
  >{\centering\arraybackslash}p{(\linewidth - 8\tabcolsep) * \real{0.2083}}
  >{\centering\arraybackslash}p{(\linewidth - 8\tabcolsep) * \real{0.2083}}
  >{\centering\arraybackslash}p{(\linewidth - 8\tabcolsep) * \real{0.2083}}
  >{\raggedright\arraybackslash}p{(\linewidth - 8\tabcolsep) * \real{0.1667}}@{}}
\toprule\noalign{}
\begin{minipage}[b]{\linewidth}\centering
Vuelta
\end{minipage} & \begin{minipage}[b]{\linewidth}\centering
\texttt{numero}
\end{minipage} & \begin{minipage}[b]{\linewidth}\centering
\texttt{cantidad\_pares}
\end{minipage} & \begin{minipage}[b]{\linewidth}\centering
Condición (\texttt{\textless{}=\ 10})
\end{minipage} & \begin{minipage}[b]{\linewidth}\raggedright
Pantalla
\end{minipage} \\
\midrule\noalign{}
\endhead
\bottomrule\noalign{}
\endlastfoot
Inicio & 2 & 0 & - & - \\
1 & 2 & 0 & \textbf{SÍ (2 \textless= 10)} & ``Encontré el par: 2'' \\
1 (Fin) & \textbf{4} & \textbf{1} & - & - \\
\emph{(\ldots)} & \emph{\ldots{}} & \emph{\ldots{}} & \emph{\ldots{}} &
\emph{``Encontré el par: 4, 6, 8''} \\
Última & 10 & 4 & \textbf{SÍ (10 \textless= 10)} & ``Encontré el par:
10'' \\
Ú. (Fin) & \textbf{12} & \textbf{5} & - & - \\
Salida & 12 & 5 & \textbf{NO (12 \textless= 10)} & ``Total de pares:
5'' \\
\end{longtable}
}

\subsubsection{Ejercicio MIENTRAS 3: Lectura con valor centinela
(``SALIR'')}\label{ejercicio-mientras-3-lectura-con-valor-centinela-salir}

\textbf{Descripción:} Recibir palabras de forma continua hasta que el
usuario ingrese el término \texttt{"SALIR"}, momento en el cual se
interrumpe la repetición.

\textbf{Pseudocódigo:}

\begin{verbatim}
INICIO
    LEER palabra
    MIENTRAS palabra != "SALIR" HACER
        IMPRIMIR "Procesando: " + palabra
        LEER palabra
    FIN MIENTRAS
    IMPRIMIR "Proceso finalizado"
FIN
\end{verbatim}

\textbf{Prueba de escritorio (ingreso de ``Hola'' y posteriormente
``SALIR''):}

{\def\LTcaptype{none} % do not increment counter
\begin{longtable}[]{@{}
  >{\centering\arraybackslash}p{(\linewidth - 6\tabcolsep) * \real{0.2632}}
  >{\centering\arraybackslash}p{(\linewidth - 6\tabcolsep) * \real{0.2632}}
  >{\centering\arraybackslash}p{(\linewidth - 6\tabcolsep) * \real{0.2632}}
  >{\raggedright\arraybackslash}p{(\linewidth - 6\tabcolsep) * \real{0.2105}}@{}}
\toprule\noalign{}
\begin{minipage}[b]{\linewidth}\centering
Vuelta
\end{minipage} & \begin{minipage}[b]{\linewidth}\centering
\texttt{palabra}
\end{minipage} & \begin{minipage}[b]{\linewidth}\centering
Condición (\texttt{!=\ "SALIR"})
\end{minipage} & \begin{minipage}[b]{\linewidth}\raggedright
Pantalla
\end{minipage} \\
\midrule\noalign{}
\endhead
\bottomrule\noalign{}
\endlastfoot
Inicio & \textbf{``Hola''} & - & \emph{Esperando texto} \\
1 & ``Hola'' & \textbf{SÍ (``Hola'' != ``SALIR'')} & ``Procesando:
Hola'' \\
1 (Fin) & \textbf{``SALIR''} & - & \emph{Esperando texto} \\
2 & ``SALIR'' & \textbf{NO (``SALIR'' != ``SALIR'' es Falso)} &
\emph{Sale del ciclo} \\
Salida & ``SALIR'' & - & ``Proceso finalizado'' \\
\end{longtable}
}

\begin{center}\rule{0.5\linewidth}{0.5pt}\end{center}

\subsection{5.2 El ciclo HACER MIENTRAS (bucle
Do-While)}\label{el-ciclo-hacer-mientras-bucle-do-while}

Evalúa la condición al \textbf{final} del bloque: -
\textbf{Comportamiento:} Ejecuta las instrucciones al menos una vez
antes de verificar si debe repetir el ciclo. - \textbf{Caso de uso:}
Validación de datos de entrada y despliegue recurrente de menús
interactivos.

\subsubsection{Ejercicio HACER MIENTRAS 1: Validación de calificación en
rango (0 a
10)}\label{ejercicio-hacer-mientras-1-validaciuxf3n-de-calificaciuxf3n-en-rango-0-a-10}

\textbf{Descripción:} Solicitar una calificación numérica y volver a
pedirla mientras el valor ingresado se encuentre fuera del rango
permitido (0 a 10).

\textbf{Pseudocódigo:}

\begin{verbatim}
INICIO
    HACER
        IMPRIMIR "Ingresa una calificación (0 a 10):"
        LEER calificacion
    MIENTRAS calificacion < 0 O calificacion > 10
    IMPRIMIR "Calificación válida guardada"
FIN
\end{verbatim}

\begin{center}
\begin{adjustbox}{max width=\textwidth,max totalheight=.72\textheight,keepaspectratio}
\begin{tikzpicture}[flujo,node distance=5mm and 14mm]
  \node[iniciofin] (n0) {INICIO};
  \node[documento,below=of n0] (n1) {"Ingresa calificación (0 a 10)"};
  \node[entradaSalida,below=of n1] (n2) {calificación};
  \node[decision,below=of n2,text width=45mm] (n3) {calificación < 0 O calificación > 10?};
  \node[documento,below=of n3] (n4) {"Calificación válida"};
  \node[iniciofin,below=of n4] (n5) {FIN};
  \draw[draw=AzulLinea] (n0) -- (n1);
  \draw[draw=TurquesaLinea] (n1) -- (n2);
  \draw[draw=RosaLinea] (n2) -- (n3);
  \draw[draw=VerdeLinea] (n3.west) --  node[etiqueta,above=2pt,pos=.45] {YES (Error, repetir)} ++(-43mm,0) |- ($(n0.south)!.5!(n1.north)$);
  \draw[draw=NaranjaLinea] (n3) -- node[etiqueta,right=2pt,pos=.45] {NO (Correcto, salir)} (n4);
  \draw[draw=TurquesaLinea] (n4) -- (n5);
\end{tikzpicture}
\end{adjustbox}
\end{center}


\textbf{Prueba de escritorio (entradas: -5, luego 15, y finalmente 8):}

{\def\LTcaptype{none} % do not increment counter
\begin{longtable}[]{@{}
  >{\centering\arraybackslash}p{(\linewidth - 8\tabcolsep) * \real{0.2083}}
  >{\centering\arraybackslash}p{(\linewidth - 8\tabcolsep) * \real{0.2083}}
  >{\centering\arraybackslash}p{(\linewidth - 8\tabcolsep) * \real{0.2083}}
  >{\centering\arraybackslash}p{(\linewidth - 8\tabcolsep) * \real{0.2083}}
  >{\raggedright\arraybackslash}p{(\linewidth - 8\tabcolsep) * \real{0.1667}}@{}}
\toprule\noalign{}
\begin{minipage}[b]{\linewidth}\centering
Vuelta
\end{minipage} & \begin{minipage}[b]{\linewidth}\centering
Acción
\end{minipage} & \begin{minipage}[b]{\linewidth}\centering
\texttt{calificacion}
\end{minipage} & \begin{minipage}[b]{\linewidth}\centering
Condición (\texttt{\textless{}\ 0\ O\ \textgreater{}\ 10})
\end{minipage} & \begin{minipage}[b]{\linewidth}\raggedright
Pantalla
\end{minipage} \\
\midrule\noalign{}
\endhead
\bottomrule\noalign{}
\endlastfoot
1 & Leer & \textbf{-5} & \textbf{SÍ (-5 \textless{} 0 es Verdadero)} &
``Ingresa calificación\ldots{}'' \\
2 & Leer & \textbf{15} & \textbf{SÍ (15 \textgreater{} 10 es Verdadero)}
& ``Ingresa calificación\ldots{}'' \\
3 & Leer & \textbf{8} & \textbf{NO (Falso, valor válido)} & ``Ingresa
calificación\ldots{}'' \\
Salida & Imprimir & 8 & - & ``Calificación válida'' \\
\end{longtable}
}

\subsubsection{Ejercicio HACER MIENTRAS 2: Menú interactivo de
opciones}\label{ejercicio-hacer-mientras-2-menuxfa-interactivo-de-opciones}

\textbf{Descripción:} Desplegar un menú con tres opciones y reiterar la
presentación hasta que el usuario elija la opción de salida (3).

\textbf{Pseudocódigo:}

\begin{verbatim}
INICIO
    HACER
        IMPRIMIR "1. Sumar"
        IMPRIMIR "2. Restar"
        IMPRIMIR "3. Salir"
        LEER opcion
    MIENTRAS opcion != 3
    IMPRIMIR "Apagando sistema"
FIN
\end{verbatim}

\begin{center}
\begin{adjustbox}{max width=\textwidth,max totalheight=.72\textheight,keepaspectratio}
\begin{tikzpicture}[flujo,node distance=5mm and 14mm]
  \node[iniciofin] (n0) {INICIO};
  \node[documento,below=of n0] (n1) {opciones de menú};
  \node[entradaSalida,below=of n1] (n2) {opción};
  \node[decision,below=of n2] (n3) {opción != 3?};
  \node[documento,below=of n3] (n4) {"Apagando sistema"};
  \node[iniciofin,below=of n4] (n5) {FIN};
  \draw[draw=AzulLinea] (n0) -- (n1);
  \draw[draw=TurquesaLinea] (n1) -- (n2);
  \draw[draw=RosaLinea] (n2) -- (n3);
  \draw[draw=VerdeLinea] (n3.west) --  node[etiqueta,above=2pt,pos=.45] {YES (Aún no quiere salir)} ++(-43mm,0) |- ($(n0.south)!.5!(n1.north)$);
  \draw[draw=NaranjaLinea] (n3) -- node[etiqueta,right=2pt,pos=.45] {NO (Eligió 3)} (n4);
  \draw[draw=TurquesaLinea] (n4) -- (n5);
\end{tikzpicture}
\end{adjustbox}
\end{center}


\textbf{Prueba de escritorio (selección de opción 1 y luego 3):}

{\def\LTcaptype{none} % do not increment counter
\begin{longtable}[]{@{}
  >{\centering\arraybackslash}p{(\linewidth - 8\tabcolsep) * \real{0.2174}}
  >{\raggedright\arraybackslash}p{(\linewidth - 8\tabcolsep) * \real{0.1739}}
  >{\centering\arraybackslash}p{(\linewidth - 8\tabcolsep) * \real{0.2174}}
  >{\centering\arraybackslash}p{(\linewidth - 8\tabcolsep) * \real{0.2174}}
  >{\raggedright\arraybackslash}p{(\linewidth - 8\tabcolsep) * \real{0.1739}}@{}}
\toprule\noalign{}
\begin{minipage}[b]{\linewidth}\centering
Vuelta
\end{minipage} & \begin{minipage}[b]{\linewidth}\raggedright
Pantalla (Menú)
\end{minipage} & \begin{minipage}[b]{\linewidth}\centering
\texttt{opcion}
\end{minipage} & \begin{minipage}[b]{\linewidth}\centering
Condición (\texttt{!=\ 3})
\end{minipage} & \begin{minipage}[b]{\linewidth}\raggedright
Acción lógica
\end{minipage} \\
\midrule\noalign{}
\endhead
\bottomrule\noalign{}
\endlastfoot
1 & ``1. Sumar, 2. Restar, 3. Salir'' & \textbf{1} & \textbf{SÍ (1 !=
3)} & \emph{Retorna a mostrar el menú} \\
2 & ``1. Sumar, 2. Restar, 3. Salir'' & \textbf{3} & \textbf{NO (3 != 3
es Falso)} & \emph{Termina el ciclo} \\
Salida & ``Apagando sistema'' & 3 & - & \emph{Finaliza} \\
\end{longtable}
}

\subsubsection{Ejercicio HACER MIENTRAS 3: Adivinar un número
secreto}\label{ejercicio-hacer-mientras-3-adivinar-un-nuxfamero-secreto}

\textbf{Descripción:} Solicitar intentos numéricos al usuario de forma
continua hasta que coincida con el número secreto predefinido (7).

\textbf{Pseudocódigo:}

\begin{verbatim}
INICIO
    secreto = 7
    HACER
        IMPRIMIR "Adivina el número secreto:"
        LEER intento
    MIENTRAS intento != secreto
    IMPRIMIR "¡Felicidades, ganaste!"
FIN
\end{verbatim}

\textbf{Prueba de escritorio (intentos: 5 y luego 7):}

{\def\LTcaptype{none} % do not increment counter
\begin{longtable}[]{@{}cclcc@{}}
\toprule\noalign{}
Vuelta & \texttt{secreto} & Pantalla & \texttt{intento} & Condición
(\texttt{!=\ secreto}) \\
\midrule\noalign{}
\endhead
\bottomrule\noalign{}
\endlastfoot
Inicio & 7 & - & - & - \\
1 & 7 & ``Adivina el número:'' & \textbf{5} & \textbf{SÍ (5 != 7)} \\
2 & 7 & ``Adivina el número:'' & \textbf{7} & \textbf{NO (7 != 7 es
Falso)} \\
Salida & 7 & ``¡Felicidades, ganaste!'' & 7 & - \\
\end{longtable}
}

\begin{center}\rule{0.5\linewidth}{0.5pt}\end{center}

\subsection{5.3 El ciclo PARA (bucle
For)}\label{el-ciclo-para-bucle-for}

Se utiliza cuando se conoce con exactitud el número de iteraciones
requeridas, integrando la inicialización, la condición de parada y el
incremento en una sola estructura de control.

\begin{itemize}
\tightlist
\item
  \textbf{Símbolo en DFD:} Hexágono alargado o bloque de control
  \texttt{\{\{\ \}\}} que especifica la variable iteradora, el rango y
  el paso.
\item
  \textbf{Sintaxis en pseudocódigo:}
  \texttt{PARA\ variable\ =\ inicio\ HASTA\ fin\ CON\ PASO\ valor\ HACER}.
\end{itemize}

\subsubsection{Ejercicio FOR 1: Conteo ascendente del 1 al
10}\label{ejercicio-for-1-conteo-ascendente-del-1-al-10}

\textbf{Descripción:} Imprimir en orden progresivo los números enteros
del 1 al 10 utilizando un paso de incremento unitario.

\textbf{Pseudocódigo:}

\begin{verbatim}
INICIO
    PARA i = 1 HASTA 10 CON PASO 1 HACER
        IMPRIMIR "El número actual es: " + i
    FIN PARA
FIN
\end{verbatim}

\textbf{Diagrama de flujo:}

\begin{center}
\begin{adjustbox}{max width=\textwidth,max totalheight=.72\textheight,keepaspectratio}
\begin{tikzpicture}[flujo,node distance=5mm and 14mm]
  \node[iniciofin] (n0) {INICIO};
  \node[proceso,below=of n0] (n1) {FOR i = 1 TO 10 STEP 1};
  \node[documento,below=of n1] (n2) {"El número actual es: " + i};
  \node[proceso,below=of n2] (n3) {Siguiente i};
  \node[iniciofin,below=of n3] (n4) {FIN};
  \draw[draw=AzulLinea] (n0) -- (n1);
  \draw[draw=TurquesaLinea] (n1) -- (n2);
  \draw[draw=RosaLinea] (n2) -- (n3);
  \draw[draw=MoradoLinea] (n3.west) --  ++(-43mm,0) |- ($(n0.south)!.5!(n1.north)$);
  \draw[draw=MoradoLinea,dashed] (n1.east) --  node[etiqueta,above=2pt,pos=.45] {Límite alcanzado} ++(43mm,0) |-  (n4.east);
\end{tikzpicture}
\end{adjustbox}
\end{center}


\textbf{Prueba de escritorio:}

{\def\LTcaptype{none} % do not increment counter
\begin{longtable}[]{@{}cccl@{}}
\toprule\noalign{}
Vuelta & Variable \texttt{i} (Automática) & Condición implícita &
Pantalla \\
\midrule\noalign{}
\endhead
\bottomrule\noalign{}
\endlastfoot
1 & \textbf{1} & SÍ (1 \textless= 10) & ``El número actual es: 1'' \\
2 & \textbf{2} & SÍ (2 \textless= 10) & ``El número actual es: 2'' \\
\emph{(\ldots)} & \emph{\ldots{}} & \emph{\ldots{}} & \emph{\ldots{}} \\
10 & \textbf{10} & SÍ (10 \textless= 10) & ``El número actual es:
10'' \\
Salida & 11 & NO (11 \textless= 10 es Falso) & \emph{Fin del
programa} \\
\end{longtable}
}

\subsubsection{Ejercicio FOR 2: Cuenta regresiva
descendente}\label{ejercicio-for-2-cuenta-regresiva-descendente}

\textbf{Descripción:} Generar una cuenta regresiva del 10 al 1 aplicando
un paso negativo de \texttt{-1}, finalizando con un mensaje de despegue.

\textbf{Pseudocódigo:}

\begin{verbatim}
INICIO
    PARA i = 10 HASTA 1 CON PASO -1 HACER
        IMPRIMIR "Lanzamiento en... " + i
    FIN PARA
    IMPRIMIR "¡Despegue!"
FIN
\end{verbatim}

\textbf{Diagrama de flujo:}

\begin{center}
\begin{adjustbox}{max width=\textwidth,max totalheight=.72\textheight,keepaspectratio}
\begin{tikzpicture}[flujo,node distance=5mm and 14mm]
  \node[iniciofin] (n0) {INICIO};
  \node[proceso,below=of n0] (n1) {FOR i = 10 TO 1 STEP -1};
  \node[documento,below=of n1] (n2) {"Lanzamiento en... " + i};
  \node[proceso,below=of n2] (n3) {Siguiente i};
  \node[documento,below=of n3] (n4) {"¡Despegue!"};
  \node[iniciofin,below=of n4] (n5) {FIN};
  \draw[draw=AzulLinea] (n0) -- (n1);
  \draw[draw=TurquesaLinea] (n1) -- (n2);
  \draw[draw=RosaLinea] (n2) -- (n3);
  \draw[draw=MoradoLinea] (n3.west) --  ++(-43mm,0) |- ($(n0.south)!.5!(n1.north)$);
  \draw[draw=MoradoLinea,dashed] (n1.east) --  node[etiqueta,above=2pt,pos=.45] {Límite alcanzado} ++(43mm,0) |-  (n4.east);
  \draw[draw=TurquesaLinea] (n4) -- (n5);
\end{tikzpicture}
\end{adjustbox}
\end{center}


\textbf{Prueba de escritorio:}

{\def\LTcaptype{none} % do not increment counter
\begin{longtable}[]{@{}cccl@{}}
\toprule\noalign{}
Vuelta & Variable \texttt{i} & Condición implícita & Pantalla \\
\midrule\noalign{}
\endhead
\bottomrule\noalign{}
\endlastfoot
1 & \textbf{10} & SÍ (10 \textgreater= 1) & ``Lanzamiento en\ldots{}
10'' \\
2 & \textbf{9} & SÍ (9 \textgreater= 1) & ``Lanzamiento en\ldots{}
9'' \\
\emph{(\ldots)} & \emph{\ldots{}} & \emph{\ldots{}} & \emph{\ldots{}} \\
10 & \textbf{1} & SÍ (1 \textgreater= 1) & ``Lanzamiento en\ldots{}
1'' \\
Salida & 0 & NO (0 \textgreater= 1 es Falso) & ``¡Despegue!'' \\
\end{longtable}
}

\subsubsection{Ejercicio FOR 3: Múltiplos de 5 con paso
modificado}\label{ejercicio-for-3-muxfaltiplos-de-5-con-paso-modificado}

\textbf{Descripción:} Imprimir la serie numérica de múltiplos de 5 desde
el 5 hasta el 50 mediante un incremento de 5 en cada iteración.

\textbf{Pseudocódigo:}

\begin{verbatim}
INICIO
    PARA iterador = 5 HASTA 50 CON PASO 5 HACER
        IMPRIMIR "Valor: " + iterador
    FIN PARA
FIN
\end{verbatim}

\textbf{Diagrama de flujo:}

\begin{center}
\begin{adjustbox}{max width=\textwidth,max totalheight=.72\textheight,keepaspectratio}
\begin{tikzpicture}[flujo,node distance=5mm and 14mm]
  \node[iniciofin] (n0) {INICIO};
  \node[proceso,below=of n0] (n1) {FOR iterador = 5 TO 50 STEP 5};
  \node[documento,below=of n1] (n2) {"Valor: " + iterador};
  \node[proceso,below=of n2] (n3) {Siguiente iterador};
  \node[iniciofin,below=of n3] (n4) {FIN};
  \draw[draw=AzulLinea] (n0) -- (n1);
  \draw[draw=TurquesaLinea] (n1) -- (n2);
  \draw[draw=RosaLinea] (n2) -- (n3);
  \draw[draw=MoradoLinea] (n3.west) --  ++(-43mm,0) |- ($(n0.south)!.5!(n1.north)$);
  \draw[draw=MoradoLinea,dashed] (n1.east) --  node[etiqueta,above=2pt,pos=.45] {Límite alcanzado} ++(43mm,0) |-  (n4.east);
\end{tikzpicture}
\end{adjustbox}
\end{center}


\textbf{Prueba de escritorio:}

{\def\LTcaptype{none} % do not increment counter
\begin{longtable}[]{@{}ccccl@{}}
\toprule\noalign{}
Vuelta & \texttt{iterador} & Salto (+5) & Condición
(\texttt{\textless{}=\ 50}) & Pantalla \\
\midrule\noalign{}
\endhead
\bottomrule\noalign{}
\endlastfoot
1 & \textbf{5} & +5 -\textgreater{} 10 & SÍ & ``Valor: 5'' \\
2 & \textbf{10} & +5 -\textgreater{} 15 & SÍ & ``Valor: 10'' \\
\emph{(\ldots)} & \emph{\ldots{}} & \emph{\ldots{}} & \emph{\ldots{}} &
\emph{\ldots{}} \\
10 & \textbf{50} & +5 -\textgreater{} 55 & SÍ & ``Valor: 50'' \\
Salida & 55 & - & NO & \emph{Fin del ciclo} \\
\end{longtable}
}

\subsubsection{Ejercicio FOR 4: Sumatoria acumulativa de 1 a
N}\label{ejercicio-for-4-sumatoria-acumulativa-de-1-a-n}

\textbf{Descripción:} Sumar los números enteros consecutivos desde 1
hasta un valor límite \texttt{N} definido por el usuario, empleando una
variable acumuladora.

\textbf{Pseudocódigo:}

\begin{verbatim}
INICIO
    LEER numero_limite
    suma_total = 0
    PARA i = 1 HASTA numero_limite CON PASO 1 HACER
        suma_total = suma_total + i
    FIN PARA
    IMPRIMIR "La sumatoria total es: " + suma_total
FIN
\end{verbatim}

\textbf{Diagrama de flujo:}

\begin{center}
\begin{adjustbox}{max width=\textwidth,max totalheight=.72\textheight,keepaspectratio}
\begin{tikzpicture}[flujo,node distance=5mm and 14mm]
  \node[iniciofin] (n0) {INICIO};
  \node[entradaSalida,below=of n0] (n1) {numero\_limite};
  \node[proceso,below=of n1] (n2) {suma\_total = 0};
  \node[proceso,below=of n2] (n3) {FOR i = 1 TO numero\_limite STEP 1};
  \node[proceso,below=of n3] (n4) {suma\_total = suma\_total + i};
  \node[proceso,below=of n4] (n5) {Siguiente i};
  \node[documento,below=of n5,text width=45mm] (n6) {"La sumatoria total es: " + suma\_total};
  \node[iniciofin,below=of n6] (n7) {FIN};
  \draw[draw=AzulLinea] (n0) -- (n1);
  \draw[draw=TurquesaLinea] (n1) -- (n2);
  \draw[draw=RosaLinea] (n2) -- (n3);
  \draw[draw=MoradoLinea] (n3) -- (n4);
  \draw[draw=AzulLinea] (n4) -- (n5);
  \draw[draw=MoradoLinea] (n5.west) --  ++(-43mm,0) |- ($(n2.south)!.5!(n3.north)$);
  \draw[draw=MoradoLinea,dashed] (n3.east) --  node[etiqueta,above=2pt,pos=.45] {Límite alcanzado} ++(43mm,0) |-  (n6.east);
  \draw[draw=MoradoLinea] (n6) -- (n7);
\end{tikzpicture}
\end{adjustbox}
\end{center}


\textbf{Prueba de escritorio (sumando hasta N = 3):}

{\def\LTcaptype{none} % do not increment counter
\begin{longtable}[]{@{}cccccl@{}}
\toprule\noalign{}
Vuelta & \texttt{limite} & \texttt{i} & \texttt{suma\_total} & Condición
& Pantalla \\
\midrule\noalign{}
\endhead
\bottomrule\noalign{}
\endlastfoot
Inicio & \textbf{3} & - & 0 & - & \emph{Esperando N} \\
1 & 3 & \textbf{1} & \textbf{1} (0+1) & SÍ (1\textless=3) &
\emph{Operación en memoria} \\
2 & 3 & \textbf{2} & \textbf{3} (1+2) & SÍ (2\textless=3) &
\emph{Operación en memoria} \\
3 & 3 & \textbf{3} & \textbf{6} (3+3) & SÍ (3\textless=3) &
\emph{Operación en memoria} \\
Salida & 3 & 4 & 6 & NO (4\textless=3) & ``La sumatoria total es: 6'' \\
\end{longtable}
}

\subsubsection{Ejercicio FOR 5: Búsqueda del valor máximo en
lote}\label{ejercicio-for-5-buxfasqueda-del-valor-muxe1ximo-en-lote}

\textbf{Descripción:} Leer 5 registros de temperatura y determinar cuál
fue el valor máximo observado, actualizando la variable \texttt{maximo}
mediante una comparación en cada iteración.

\textbf{Pseudocódigo:}

\begin{verbatim}
INICIO
    maximo = -9999
    PARA dia = 1 HASTA 5 CON PASO 1 HACER
        LEER temperatura
        SI temperatura > maximo ENTONCES
            maximo = temperatura
        FIN SI
    FIN PARA
    IMPRIMIR "La temperatura máxima de la semana fue: " + maximo
FIN
\end{verbatim}

\textbf{Diagrama de flujo:}

\begin{center}
\begin{adjustbox}{max width=\textwidth,max totalheight=.72\textheight,keepaspectratio}
\begin{tikzpicture}[flujo,node distance=5mm and 14mm]
  \node[iniciofin] (n0) {INICIO};
  \node[proceso,below=of n0] (n1) {máximo = -9999};
  \node[proceso,below=of n1] (n2) {FOR dia = 1 TO 5 STEP 1};
  \node[entradaSalida,below=of n2] (n3) {temperatura};
  \node[decision,below=of n3] (n4) {temperatura > máximo?};
  \node[proceso,below=of n4] (n5) {máximo = temperatura};
  \node[proceso,below=of n5] (n6) {Siguiente dia};
  \node[documento,below=of n6] (n7) {"La temperatura máxima... " + máximo};
  \node[iniciofin,below=of n7] (n8) {FIN};
  \draw[draw=AzulLinea] (n0) -- (n1);
  \draw[draw=TurquesaLinea] (n1) -- (n2);
  \draw[draw=RosaLinea] (n2) -- (n3);
  \draw[draw=MoradoLinea] (n3) -- (n4);
  \draw[draw=VerdeLinea] (n4) -- node[etiqueta,right=2pt,pos=.45] {YES} (n5);
  \draw[draw=NaranjaLinea] (n4.east) --  node[etiqueta,above=2pt,pos=.45] {NO} ++(43mm,0) |-  (n6.east);
  \draw[draw=RosaLinea] (n5) -- (n6);
  \draw[draw=MoradoLinea] (n6.west) --  ++(-43mm,0) |- ($(n1.south)!.5!(n2.north)$);
  \draw[draw=MoradoLinea,dashed] (n2.east) --  node[etiqueta,above=2pt,pos=.45] {Límite alcanzado} ++(48mm,0) |-  (n7.east);
  \draw[draw=TurquesaLinea] (n7) -- (n8);
\end{tikzpicture}
\end{adjustbox}
\end{center}


\textbf{Prueba de escritorio (temperaturas: 25, 28, 22, 30, 29):}

{\def\LTcaptype{none} % do not increment counter
\begin{longtable}[]{@{}
  >{\centering\arraybackslash}p{(\linewidth - 10\tabcolsep) * \real{0.1667}}
  >{\centering\arraybackslash}p{(\linewidth - 10\tabcolsep) * \real{0.1667}}
  >{\centering\arraybackslash}p{(\linewidth - 10\tabcolsep) * \real{0.1667}}
  >{\centering\arraybackslash}p{(\linewidth - 10\tabcolsep) * \real{0.1667}}
  >{\centering\arraybackslash}p{(\linewidth - 10\tabcolsep) * \real{0.1667}}
  >{\centering\arraybackslash}p{(\linewidth - 10\tabcolsep) * \real{0.1667}}@{}}
\toprule\noalign{}
\begin{minipage}[b]{\linewidth}\centering
Vuelta
\end{minipage} & \begin{minipage}[b]{\linewidth}\centering
\texttt{dia}
\end{minipage} & \begin{minipage}[b]{\linewidth}\centering
\texttt{maximo}
\end{minipage} & \begin{minipage}[b]{\linewidth}\centering
\texttt{temperatura}
\end{minipage} & \begin{minipage}[b]{\linewidth}\centering
¿Es \textgreater{} máximo?
\end{minipage} & \begin{minipage}[b]{\linewidth}\centering
Nuevo máximo
\end{minipage} \\
\midrule\noalign{}
\endhead
\bottomrule\noalign{}
\endlastfoot
Inicio & - & -9999 & - & - & - \\
1 & 1 & -9999 & \textbf{25} & SÍ (25 \textgreater{} -9999) &
\textbf{25} \\
2 & 2 & 25 & \textbf{28} & SÍ (28 \textgreater{} 25) & \textbf{28} \\
3 & 3 & 28 & \textbf{22} & NO (22 \textless= 28) & 28 \\
4 & 4 & 28 & \textbf{30} & SÍ (30 \textgreater{} 28) & \textbf{30} \\
5 & 5 & 30 & \textbf{29} & NO (29 \textless= 30) & 30 \\
Salida & 6 & 30 & - & (Fin) & ``La temperatura máxima\ldots{} 30'' \\
\end{longtable}
}

\subsubsection{Ejercicio FOR 6: Tabla de
multiplicar}\label{ejercicio-for-6-tabla-de-multiplicar}

\textbf{Descripción:} Solicitar un número entero y mostrar su tabla de
multiplicar del 1 al 10, calculando y formateando el producto en cada
paso.

\textbf{Pseudocódigo:}

\begin{verbatim}
INICIO
    LEER tabla_del
    PARA multiplicador = 1 HASTA 10 CON PASO 1 HACER
        resultado = tabla_del * multiplicador
        IMPRIMIR tabla_del + " x " + multiplicador + " = " + resultado
    FIN PARA
FIN
\end{verbatim}

\textbf{Diagrama de flujo:}

\begin{center}
\begin{adjustbox}{max width=\textwidth,max totalheight=.72\textheight,keepaspectratio}
\begin{tikzpicture}[flujo,node distance=5mm and 14mm]
  \node[iniciofin] (n0) {INICIO};
  \node[entradaSalida,below=of n0] (n1) {tabla\_del};
  \node[proceso,below=of n1] (n2) {FOR multiplicador = 1 TO 10 STEP 1};
  \node[proceso,below=of n2,text width=45mm] (n3) {resultado = tabla\_del * multiplicador};
  \node[documento,below=of n3,text width=45mm] (n4) {tabla\_del + " x " + multiplicador + " = " + resultado};
  \node[proceso,below=of n4] (n5) {Siguiente multiplicador};
  \node[iniciofin,below=of n5] (n6) {FIN};
  \draw[draw=AzulLinea] (n0) -- (n1);
  \draw[draw=TurquesaLinea] (n1) -- (n2);
  \draw[draw=RosaLinea] (n2) -- (n3);
  \draw[draw=MoradoLinea] (n3) -- (n4);
  \draw[draw=AzulLinea] (n4) -- (n5);
  \draw[draw=MoradoLinea] (n5.west) --  ++(-43mm,0) |- ($(n1.south)!.5!(n2.north)$);
  \draw[draw=MoradoLinea,dashed] (n2.east) --  node[etiqueta,above=2pt,pos=.45] {Límite alcanzado} ++(43mm,0) |-  (n6.east);
\end{tikzpicture}
\end{adjustbox}
\end{center}


\textbf{Prueba de escritorio (tabla del 7):}

{\def\LTcaptype{none} % do not increment counter
\begin{longtable}[]{@{}ccccl@{}}
\toprule\noalign{}
Vuelta & \texttt{tabla\_del} & \texttt{multiplicador} &
\texttt{resultado} & Pantalla \\
\midrule\noalign{}
\endhead
\bottomrule\noalign{}
\endlastfoot
Inicio & \textbf{7} & - & - & \emph{Esperando número} \\
1 & 7 & \textbf{1} & \textbf{7} & ``7 x 1 = 7'' \\
2 & 7 & \textbf{2} & \textbf{14} & ``7 x 2 = 14'' \\
\emph{(\ldots)} & \emph{\ldots{}} & \emph{\ldots{}} & \emph{\ldots{}} &
\emph{\ldots{}} \\
10 & 7 & \textbf{10} & \textbf{70} & ``7 x 10 = 70'' \\
Salida & 7 & 11 & - & \emph{Fin del ciclo} \\
\end{longtable}
}

\begin{center}\rule{0.5\linewidth}{0.5pt}\end{center}

\subsection{5.4 Ejercicios combinados con múltiples
estructuras}\label{ejercicios-combinados-con-muxfaltiples-estructuras}

En aplicaciones reales es común secuenciar o anidar distintos tipos de
ciclos y condicionales según los requerimientos del problema.

\subsubsection{Ejercicio Combinado 1: Validación previa con lectura
múltiple (Hacer Mientras +
Mientras)}\label{ejercicio-combinado-1-validaciuxf3n-previa-con-lectura-muxfaltiple-hacer-mientras-mientras}

\textbf{Descripción:} Validar que la cantidad de alumnos a procesar sea
estrictamente mayor a cero y, posteriormente, solicitar la calificación
individual de cada uno mediante un ciclo de control.

\textbf{Pseudocódigo:}

\begin{verbatim}
INICIO
    HACER
        IMPRIMIR "Ingresa la cantidad de alumnos (mayor a 0):"
        LEER n_alumnos
    MIENTRAS n_alumnos <= 0
    
    contador = 1
    MIENTRAS contador <= n_alumnos HACER
        IMPRIMIR "Ingresa la calificación del alumno " + contador
        LEER calif
        contador = contador + 1
    FIN MIENTRAS
    IMPRIMIR "Todas las calificaciones registradas."
FIN
\end{verbatim}

\textbf{Prueba de escritorio (entradas: -2, luego 2, y calificaciones 8
y 9):}

{\def\LTcaptype{none} % do not increment counter
\begin{longtable}[]{@{}
  >{\centering\arraybackslash}p{(\linewidth - 6\tabcolsep) * \real{0.2632}}
  >{\centering\arraybackslash}p{(\linewidth - 6\tabcolsep) * \real{0.2632}}
  >{\centering\arraybackslash}p{(\linewidth - 6\tabcolsep) * \real{0.2632}}
  >{\raggedright\arraybackslash}p{(\linewidth - 6\tabcolsep) * \real{0.2105}}@{}}
\toprule\noalign{}
\begin{minipage}[b]{\linewidth}\centering
Paso
\end{minipage} & \begin{minipage}[b]{\linewidth}\centering
Variable / Acción
\end{minipage} & \begin{minipage}[b]{\linewidth}\centering
Valor / Condición
\end{minipage} & \begin{minipage}[b]{\linewidth}\raggedright
Pantalla
\end{minipage} \\
\midrule\noalign{}
\endhead
\bottomrule\noalign{}
\endlastfoot
Val 1 & \texttt{n\_alumnos} & \textbf{-2} (Inválido: -2 \textless= 0) &
``Ingresa cantidad de alumnos\ldots{}'' \\
Val 2 & \texttt{n\_alumnos} & \textbf{2} (Válido: 2 \textless= 0 es
Falso) & ``Ingresa cantidad de alumnos\ldots{}'' \\
Ciclo 1 & \texttt{contador} & 1 & ``Ingresa la calificación del alumno
1'' \\
Ciclo 1 & \texttt{calif} & \textbf{8} & - \\
Ciclo 2 & \texttt{contador} & 2 & ``Ingresa la calificación del alumno
2'' \\
Ciclo 2 & \texttt{calif} & \textbf{9} & - \\
Fin & Salida & - & ``Todas las calificaciones registradas.'' \\
\end{longtable}
}

\subsubsection{Ejercicio Combinado 2: Bucle interactivo con fases fijas
(Mientras +
Para)}\label{ejercicio-combinado-2-bucle-interactivo-con-fases-fijas-mientras-para}

\textbf{Descripción:} Recibir palabras de forma continua hasta que se
ingrese \texttt{"FIN"}. Por cada palabra procesada, ejecutar un ciclo
interno fijo de 3 pasos que emite asteriscos de confirmación.

\textbf{Pseudocódigo:}

\begin{verbatim}
INICIO
    LEER palabra
    MIENTRAS palabra != "FIN" HACER
        IMPRIMIR "Analizando: " + palabra
        PARA fase = 1 HASTA 3 CON PASO 1 HACER
            IMPRIMIR "*"
        FIN PARA
        IMPRIMIR "Análisis completado."
        LEER palabra
    FIN MIENTRAS
FIN
\end{verbatim}

\textbf{Prueba de escritorio (ingreso de ``Hola'' y luego ``FIN''):}

{\def\LTcaptype{none} % do not increment counter
\begin{longtable}[]{@{}
  >{\centering\arraybackslash}p{(\linewidth - 6\tabcolsep) * \real{0.2632}}
  >{\centering\arraybackslash}p{(\linewidth - 6\tabcolsep) * \real{0.2632}}
  >{\centering\arraybackslash}p{(\linewidth - 6\tabcolsep) * \real{0.2632}}
  >{\raggedright\arraybackslash}p{(\linewidth - 6\tabcolsep) * \real{0.2105}}@{}}
\toprule\noalign{}
\begin{minipage}[b]{\linewidth}\centering
Paso
\end{minipage} & \begin{minipage}[b]{\linewidth}\centering
Variable
\end{minipage} & \begin{minipage}[b]{\linewidth}\centering
Valor
\end{minipage} & \begin{minipage}[b]{\linewidth}\raggedright
Pantalla
\end{minipage} \\
\midrule\noalign{}
\endhead
\bottomrule\noalign{}
\endlastfoot
Mientras & \texttt{palabra} & \textbf{``Hola''} & \emph{Esperando
texto} \\
Mientras & Condición & SÍ (``Hola'' != ``FIN'') & ``Analizando:
Hola'' \\
Para & \texttt{fase} 1 & 1 & ``*'' \\
Para & \texttt{fase} 2 & 2 & ``*'' \\
Para & \texttt{fase} 3 & 3 & ``*'' \\
Para & Fin & - & ``Análisis completado.'' \\
Mientras & \texttt{palabra} & \textbf{``FIN''} & \emph{Esperando
texto} \\
Mientras & Condición & NO (Término del ciclo) & \emph{Finaliza el
programa} \\
\end{longtable}
}

\subsubsection{Ejercicio Combinado 3: Menú interactivo con secuencias
ascendente y descendente (Hacer Mientras +
Para)}\label{ejercicio-combinado-3-menuxfa-interactivo-con-secuencias-ascendente-y-descendente-hacer-mientras-para}

\textbf{Descripción:} Desplegar un menú que permite contar del 1 al 3
(opción 1), del 3 al 1 (opción 2) o salir del sistema (opción 3),
empleando condicionales anidados dentro de una estructura repetitiva.

\textbf{Pseudocódigo:}

\begin{verbatim}
INICIO
    HACER
        IMPRIMIR "1. Subir, 2. Bajar, 3. Salir"
        LEER opc
        SI opc == 1 ENTONCES
            PARA i = 1 HASTA 3 CON PASO 1 HACER
                IMPRIMIR "Subiendo... " + i
            FIN PARA
        SINO
            SI opc == 2 ENTONCES
                PARA i = 3 HASTA 1 CON PASO -1 HACER
                    IMPRIMIR "Bajando... " + i
                FIN PARA
            FIN SI
        FIN SI
    MIENTRAS opc != 3
    IMPRIMIR "Apagado"
FIN
\end{verbatim}

\textbf{Prueba de escritorio (opción 1 y luego opción 3):}

{\def\LTcaptype{none} % do not increment counter
\begin{longtable}[]{@{}
  >{\centering\arraybackslash}p{(\linewidth - 6\tabcolsep) * \real{0.2632}}
  >{\centering\arraybackslash}p{(\linewidth - 6\tabcolsep) * \real{0.2632}}
  >{\centering\arraybackslash}p{(\linewidth - 6\tabcolsep) * \real{0.2632}}
  >{\raggedright\arraybackslash}p{(\linewidth - 6\tabcolsep) * \real{0.2105}}@{}}
\toprule\noalign{}
\begin{minipage}[b]{\linewidth}\centering
Paso
\end{minipage} & \begin{minipage}[b]{\linewidth}\centering
Acción / Variable
\end{minipage} & \begin{minipage}[b]{\linewidth}\centering
Valor
\end{minipage} & \begin{minipage}[b]{\linewidth}\raggedright
Pantalla
\end{minipage} \\
\midrule\noalign{}
\endhead
\bottomrule\noalign{}
\endlastfoot
Menú & \texttt{opc} & \textbf{1} & ``1. Subir, 2. Bajar, 3. Salir'' \\
Cond & Evaluar opción 1 & SÍ (1 == 1) & - \\
Para 1 & \texttt{i} (Vueltas 1, 2, 3) & 1..3 & ``Subiendo\ldots{} 1'',
``Subiendo\ldots{} 2'', ``Subiendo\ldots{} 3'' \\
Bucle & Retorno del menú & SÍ (1 != 3) & \emph{Vuelve al menú
principal} \\
Menú & \texttt{opc} & \textbf{3} & ``1. Subir, 2. Bajar, 3. Salir'' \\
Cond & Evaluar opción 3 & NO (a 1 y a 2) & - \\
Bucle & Retorno del menú & NO (3 != 3 es Falso) & ``Apagado'' \\
\end{longtable}
}

\begin{center}\rule{0.5\linewidth}{0.5pt}\end{center}

\section{6. Clasificación de errores en la
programación}\label{clasificaciuxf3n-de-errores-en-la-programaciuxf3n}

Al diseñar e implementar algoritmos se presentan distintos tipos de
errores según la etapa en la que se manifiestan y la herramienta capaz
de detectarlos:

\begin{enumerate}
\def\labelenumi{\arabic{enumi}.}
\item
  \textbf{Errores de sintaxis:} Ocurren al infringir las reglas
  gramaticales o de estructura del lenguaje (por ejemplo, omitir el
  cierre de comillas o escribir incorrectamente una palabra reservada
  como \texttt{IMPRIMR}). El compilador o intérprete los identifica
  antes de la ejecución y señala la línea donde se originan.
\item
  \textbf{Errores de ejecución (Runtime errors):} Aparecen cuando las
  instrucciones son sintácticamente válidas, pero solicitan operaciones
  no permitidas durante el funcionamiento del programa, tales como
  intentar dividir entre cero o acceder a un recurso no disponible.
\item
  \textbf{Errores lógicos:} Se presentan cuando el programa se ejecuta
  sin interrupciones pero entrega resultados incorrectos debido a una
  falla en el razonamiento, las fórmulas o el flujo de control. Las
  herramientas automáticas no pueden detectarlos directamente, por lo
  que su prevención y corrección dependen de la prueba de escritorio y
  del análisis del desarrollador.
\end{enumerate}

\subsection{6.1 Tiempo de compilación vs.~tiempo de
ejecución}\label{tiempo-de-compilaciuxf3n-vs.-tiempo-de-ejecuciuxf3n}

Para comprender la naturaleza de estos errores es útil distinguir las
dos etapas principales en el ciclo de un programa:

\begin{itemize}
\item
  \textbf{Tiempo de compilación (Compile Time):} Fase previa en la que
  el compilador analiza el código fuente, valida su sintaxis y lo
  traduce a instrucciones ejecutables. Los errores de sintaxis se
  identifican exclusivamente en este momento, impidiendo la generación
  del programa final hasta ser corregidos.
\item
  \textbf{Tiempo de ejecución (Run Time):} Fase en la que el programa
  compilado o interpretado se ejecuta en la memoria de la computadora,
  procesa entradas y genera salidas. En esta etapa se manifiestan los
  errores de ejecución y se observan las consecuencias de los errores
  lógicos.
\end{itemize}

\subsection{6.2 Errores lógicos y de control comunes en
diagramas}\label{errores-luxf3gicos-y-de-control-comunes-en-diagramas}

Durante el diseño de algoritmos iterativos suelen presentarse
situaciones específicas que conviene anticipar:

\begin{itemize}
\tightlist
\item
  \textbf{Bucle infinito:} Se produce cuando la condición de salida
  depende de una variable cuyo valor nunca cambia dentro del ciclo (por
  ejemplo, al omitir el incremento
  \texttt{contador\ =\ contador\ +\ 1}), provocando que la condición sea
  siempre verdadera.
\item
  \textbf{Ciclo inactivo:} Ocurre en estructuras \texttt{MIENTRAS}
  cuando la variable de control se inicializa con un valor que hace
  falsa la condición de entrada desde el primer momento, omitiendo el
  bloque por completo.
\item
  \textbf{Error por desplazamiento de un paso (Off-by-One):} Sucede al
  confundir operadores como \texttt{\textless{}} y
  \texttt{\textless{}=}, lo que ocasiona que el ciclo realice una
  iteración más o una menos de las previstas.
\item
  \textbf{División por cero no controlada:} Ocurre cuando se calcula un
  promedio o tasa sin validar previamente si el divisor es cero (por
  ejemplo, con \texttt{N\ =\ 0}), lo cual genera una interrupción de
  ejecución si no se añade una verificación previa.
\end{itemize}

\begin{center}\rule{0.5\linewidth}{0.5pt}\end{center}

\section{7. Ejercicio Integrador: Cajero automático
básico}\label{ejercicio-integrador-cajero-automuxe1tico-buxe1sico}

A continuación se presenta un ejercicio que combina las principales
estructuras abordadas: ciclo \texttt{HACER\ MIENTRAS} para la navegación
del menú, condicionales anidados para la validación de fondos y ciclo
\texttt{PARA} para la rutina final de desconexión.

\textbf{Planteamiento:} Diseñar un sistema de cajero automático con un
saldo inicial de 1000 unidades monetarias que cumpla con los siguientes
puntos:

\begin{enumerate}
\def\labelenumi{\arabic{enumi}.}
\tightlist
\item
  Desplegar un menú interactivo continuo con tres opciones: 1. Consultar
  saldo, 2. Realizar retiro, 3. Salir.
\item
  Al seleccionar retiro, solicitar el monto y verificar que no exceda el
  saldo disponible; si los fondos son suficientes, actualizar el saldo,
  y en caso contrario, mostrar un mensaje de fondos insuficientes.
\item
  Al elegir la opción de salida (3), ejecutar un ciclo de 3 pasos que
  simule la desconexión del servidor y emitir un mensaje de
  agradecimiento.
\end{enumerate}

\textbf{Pseudocódigo:}

\begin{verbatim}
INICIO
    saldo = 1000
    HACER
        IMPRIMIR "1. Consultar"
        IMPRIMIR "2. Retirar"
        IMPRIMIR "3. Salir"
        LEER opcion
        
        SI opcion == 1 ENTONCES
            IMPRIMIR "Tu saldo es: " + saldo
        SINO
            SI opcion == 2 ENTONCES
                IMPRIMIR "Ingresa monto a retirar:"
                LEER retiro
                SI retiro <= saldo ENTONCES
                    saldo = saldo - retiro
                    IMPRIMIR "Retiro exitoso. Nuevo saldo: " + saldo
                SINO
                    IMPRIMIR "Fondos insuficientes"
                FIN SI
            FIN SI
        FIN SI
    MIENTRAS opcion != 3
    
    IMPRIMIR "Iniciando protocolo de desconexión..."
    PARA fase = 1 HASTA 3 CON PASO 1 HACER
        IMPRIMIR "Cerrando fase de seguridad: " + fase
    FIN PARA
    IMPRIMIR "Gracias por usar nuestro banco"
FIN
\end{verbatim}

\clearpage

\textbf{Diagrama de flujo:}

\begin{center}
\begin{adjustbox}{max width=\textwidth,max totalheight=.72\textheight,keepaspectratio}
\begin{tikzpicture}[flujo,node distance=5mm and 14mm]
  \node[iniciofin] (a) {INICIO};
  \node[proceso,below=of a] (b) {saldo = 1000};
  \node[entradaSalida,below=of b] (c) {menú y opción};
  \node[decision,below=of c] (d) {opción == 1?};
  \node[documento,left=42mm of d] (e) {saldo};
  \node[decision,below=of d] (f) {opción == 2?};
  \node[entradaSalida,below=of f] (g) {retiro};
  \node[decision,below=of g] (h) {retiro <= saldo?};
  \node[proceso,text width=45mm,below=of h] (i) {saldo = saldo - retiro\\"Retiro exitoso"};
  \node[documento,right=42mm of h] (j) {"Fondos insuficientes"};
  \node[decision,below=of i] (k) {opción != 3?};
  \node[documento,below=of k] (l) {"Iniciando protocolo..."};
  \node[proceso,below=of l] (m) {FOR fase = 1 TO 3 STEP 1};
  \node[documento,below=of m] (n) {"Cerrando fase... " + fase};
  \node[proceso,below=of n] (o) {Siguiente fase};
  \node[documento,right=42mm of m] (p) {"Gracias"};
  \node[iniciofin,below=of p] (q) {FIN};
  \node[conector,left=52mm of c] (ca2) {A};
  \node[conector,left=52mm of k] (ca1) {A};
  \node[conector,right=52mm of f] (cc1) {C};
  \node[conector,right=52mm of k] (cc2) {C};
  \node[conector,left=52mm of m] (cb2) {B};
  \node[conector,left=52mm of o] (cb1) {B};
  \draw[draw=AzulLinea] (a) -- (b);
  \draw[draw=TurquesaLinea] (b) -- (c);
  \draw[draw=RosaLinea] (c) -- (d);
  \draw[draw=VerdeLinea] (d.west) -- node[etiqueta,above=2pt] {SÍ} (e.east);
  \draw[draw=MoradoLinea] (e.south) |- ($(i.south)!.25!(k.north)$);
  \draw[draw=NaranjaLinea] (d) -- node[etiqueta,right=2pt] {NO} (f);
  \draw[draw=VerdeLinea] (f) -- node[etiqueta,right=2pt] {SÍ} (g);
  \draw[draw=NaranjaLinea] (f.east) -- node[etiqueta,above=2pt] {NO} (cc1);
  \draw[draw=MoradoLinea] (cc2) |- ($(i.south)!.5!(k.north)$);
  \draw[draw=AzulLinea] (g) -- (h);
  \draw[draw=VerdeLinea] (h) -- node[etiqueta,right=2pt] {SÍ} (i);
  \draw[draw=NaranjaLinea] (h.east) -- node[etiqueta,above=2pt] {NO} (j.west);
  \draw[draw=TurquesaLinea] (i) -- (k);
  \draw[draw=MoradoLinea] (j.south) |- ($(i.south)!.75!(k.north)$);
  \draw[draw=VerdeLinea] (k.west) -- node[etiqueta,above=2pt] {SÍ} (ca1);
  \draw[draw=MoradoLinea] (ca2) |- ($(b.south)!.5!(c.north)$);
  \draw[draw=NaranjaLinea] (k) -- node[etiqueta,right=2pt] {NO} (l);
  \draw[draw=AzulLinea] (l) -- (m);
  \draw[draw=TurquesaLinea] (m) -- (n);
  \draw[draw=RosaLinea] (n) -- (o);
  \draw[draw=MoradoLinea] (o.west) -- (cb1);
  \draw[draw=MoradoLinea] (cb2) |- ($(l.south)!.5!(m.north)$);
  \draw[draw=MoradoLinea,dashed] (m.east) -- node[etiqueta,above=2pt] {Límite alcanzado} (p.west);
  \draw[draw=AzulLinea] (p) -- (q);
\end{tikzpicture}
\end{adjustbox}
\end{center}


\textbf{Prueba de escritorio (intento de retiro de 1500, retiro de 200,
consulta de saldo y salida):}

{\def\LTcaptype{none} % do not increment counter
\begin{longtable}[]{@{}
  >{\centering\arraybackslash}p{(\linewidth - 6\tabcolsep) * \real{0.2778}}
  >{\raggedright\arraybackslash}p{(\linewidth - 6\tabcolsep) * \real{0.2222}}
  >{\centering\arraybackslash}p{(\linewidth - 6\tabcolsep) * \real{0.2778}}
  >{\raggedright\arraybackslash}p{(\linewidth - 6\tabcolsep) * \real{0.2222}}@{}}
\toprule\noalign{}
\begin{minipage}[b]{\linewidth}\centering
Paso
\end{minipage} & \begin{minipage}[b]{\linewidth}\raggedright
Acción lógica / Variable
\end{minipage} & \begin{minipage}[b]{\linewidth}\centering
Valor
\end{minipage} & \begin{minipage}[b]{\linewidth}\raggedright
Pantalla (Salida)
\end{minipage} \\
\midrule\noalign{}
\endhead
\bottomrule\noalign{}
\endlastfoot
Inicio & \texttt{saldo} inicial & 1000 & - \\
\textbf{B1} & \textbf{--- Menú principal ---} & - & ``1. Consultar, 2.
Retirar, 3. Salir'' \\
Leer 1 & \texttt{opcion} & \textbf{2} & \emph{Esperando opción} \\
Cond & ¿Opción 2? & SÍ & ``Ingresa monto a retirar:'' \\
Leer & \texttt{retiro} & \textbf{1500} & \emph{Esperando monto} \\
Cond & ¿\texttt{retiro} \textless= \texttt{saldo}? (1500 \textless=
1000) & NO & ``Fondos insuficientes'' \\
Eval & ¿\texttt{opcion} != 3? & SÍ & \emph{Retorna al menú} \\
\textbf{B2} & \textbf{--- Menú principal ---} & - & ``1. Consultar, 2.
Retirar, 3. Salir'' \\
Leer 2 & \texttt{opcion} & \textbf{2} & \emph{Esperando opción} \\
Cond & ¿Opción 2? & SÍ & ``Ingresa monto a retirar:'' \\
Leer & \texttt{retiro} & \textbf{200} & \emph{Esperando monto} \\
Cond & ¿\texttt{retiro} \textless= \texttt{saldo}? (200 \textless= 1000)
& SÍ & \emph{Operación en memoria} \\
Proc & \texttt{saldo} (1000 - 200) & \textbf{800} & ``Retiro exitoso.
Nuevo saldo: 800'' \\
Eval & ¿\texttt{opcion} != 3? & SÍ & \emph{Retorna al menú} \\
\textbf{B3} & \textbf{--- Menú principal ---} & - & ``1. Consultar, 2.
Retirar, 3. Salir'' \\
Leer 3 & \texttt{opcion} & \textbf{1} & \emph{Esperando opción} \\
Cond & ¿Opción 1? & SÍ & ``Tu saldo es: 800'' \\
Eval & ¿\texttt{opcion} != 3? & SÍ & \emph{Retorna al menú} \\
\textbf{B4} & \textbf{--- Menú principal ---} & - & ``1. Consultar, 2.
Retirar, 3. Salir'' \\
Leer 4 & \texttt{opcion} & \textbf{3} & \emph{Esperando opción} \\
Eval & ¿\texttt{opcion} != 3? & NO & \emph{Concluye el menú} \\
Salida & Fin del menú & - & ``Iniciando protocolo de
desconexión\ldots{}'' \\
Para & \texttt{fase} = 1 & 1 & ``Cerrando fase de seguridad: 1'' \\
Para & \texttt{fase} = 2 & 2 & ``Cerrando fase de seguridad: 2'' \\
Para & \texttt{fase} = 3 & 3 & ``Cerrando fase de seguridad: 3'' \\
Fin & Mensaje final & - & ``Gracias por usar nuestro banco'' \\
\end{longtable}
}

\begin{center}\rule{0.5\linewidth}{0.5pt}\end{center}

\section{8. Recomendaciones y siguientes
pasos}\label{recomendaciones-y-siguientes-pasos}

El diseño de algoritmos parte de instrucciones elementales y avanza
hacia estructuras organizadas mediante reglas claras de flujo. Para
consolidar esta práctica, se sugieren las siguientes recomendaciones:

\begin{itemize}
\tightlist
\item
  \textbf{Herramientas de diagramación:} Para diagramar de forma digital
  pueden utilizarse herramientas como \textbf{yEd Graph Editor},
  \textbf{Draw.io} o \textbf{Lucidchart}, así como el dibujo tradicional
  en papel manteniendo la consistencia de los símbolos estándar.
\item
  \textbf{Verificación constante:} Aplicar la prueba de escritorio a
  cada solución permite detectar inconsistencias lógicas y evaluar casos
  límite antes de pasar a la codificación.
\item
  \textbf{Estructuras de datos colectivas:} Tras dominar variables
  simples, el siguiente paso comprende el uso de \textbf{arreglos
  (arrays)} o \textbf{listas}, los cuales permiten almacenar y procesar
  conjuntos múltiples de datos bajo un mismo identificador.
\item
  \textbf{Modularidad:} Conforme aumenta la complejidad de los
  programas, las operaciones repetitivas o extensas se organizan en
  subprocesos independientes (\textbf{funciones y métodos}), facilitando
  el mantenimiento y la reutilización del código.
\end{itemize}

\begin{center}\rule{0.5\linewidth}{0.5pt}\end{center}

\section{9. Referencias bibliográficas de
consulta}\label{referencias-bibliogruxe1ficas-de-consulta}

\begin{enumerate}
\def\labelenumi{\arabic{enumi}.}
\tightlist
\item
  \textbf{Joyanes Aguilar, L. (2008).} \emph{Fundamentos de
  programación: Algoritmos, estructuras de datos y objetos} (4.ª ed.).
  McGraw-Hill Interamericana.
\item
  \textbf{Cairó Battistutti, O. (2005).} \emph{Metodología de la
  programación: Algoritmos, diagramas de flujo y programas} (3.ª ed.).
  Alfaomega.
\item
  \textbf{Cormen, T. H., Leiserson, C. E., Rivest, R. L., \& Stein, C.
  (2009).} \emph{Introduction to Algorithms} (3.ª ed.). MIT Press.
\end{enumerate}


\end{document}
